% pandoc 模板:版式与 assignment01/handout/assignment01.tex 相同,
% 正文由 handout/src/*.md 生成,不要手工编辑生成出来的 .tex。
\documentclass[11pt,a4paper,fontset=fandol]{ctexart}

\usepackage{geometry}
\geometry{margin=2.3cm}

\usepackage{amsmath,amssymb}
\usepackage{booktabs,array,multirow,longtable}
\usepackage{calc,etoolbox}
\usepackage{enumitem}
\setlist{nosep,leftmargin=2em}
\usepackage{xcolor}
\definecolor{accent}{RGB}{28,60,120}
\definecolor{codebg}{RGB}{248,248,246}
\usepackage[most]{tcolorbox}
\usepackage{listings}
\usepackage{fvextra}
\IfFontExistsTF{DejaVu Sans Mono}{\setmonofont{DejaVu Sans Mono}[Scale=MatchLowercase]}{}
\definecolor{strgreen}{RGB}{21,101,42}
\lstset{
  basicstyle=\ttfamily\small,
  backgroundcolor=\color{codebg},
  frame=single, rulecolor=\color{gray!40}, framerule=0.4pt,
  breaklines=true, columns=fullflexible, keepspaces=true,
  showstringspaces=false,
  language=bash,
  morekeywords={make,nvcc,uv,pytest,python,cd,ncu,nvdisasm},
  commentstyle=\color{gray!130}\itshape,
  keywordstyle=\color{accent}\bfseries,
  stringstyle=\color{strgreen},
  xleftmargin=4pt, xrightmargin=4pt,
}
\usepackage{hyperref}
\hypersetup{colorlinks, linkcolor=accent, urlcolor=accent}
\emergencystretch=3em
\pagestyle{plain}

% pandoc 产物需要的宏
\providecommand{\tightlist}{\setlength{\itemsep}{0pt}\setlength{\parskip}{0pt}}
\providecommand{\passthrough}[1]{#1}
\setlength{\LTpre}{0.6em}
\setlength{\LTpost}{0.6em}
\AtBeginEnvironment{longtable}{\small\setlength{\tabcolsep}{3pt}}

\usepackage{color}
\usepackage{fancyvrb}
\newcommand{\VerbBar}{|}
\newcommand{\VERB}{\Verb[commandchars=\\\{\}]}
\DefineVerbatimEnvironment{Highlighting}{Verbatim}{commandchars=\\\{\}}
% Add ',fontsize=\small' for more characters per line
\newenvironment{Shaded}{}{}
\newcommand{\AlertTok}[1]{\textcolor[rgb]{1.00,0.00,0.00}{\textbf{#1}}}
\newcommand{\AnnotationTok}[1]{\textcolor[rgb]{0.38,0.63,0.69}{\textbf{\textit{#1}}}}
\newcommand{\AttributeTok}[1]{\textcolor[rgb]{0.49,0.56,0.16}{#1}}
\newcommand{\BaseNTok}[1]{\textcolor[rgb]{0.25,0.63,0.44}{#1}}
\newcommand{\BuiltInTok}[1]{\textcolor[rgb]{0.00,0.50,0.00}{#1}}
\newcommand{\CharTok}[1]{\textcolor[rgb]{0.25,0.44,0.63}{#1}}
\newcommand{\CommentTok}[1]{\textcolor[rgb]{0.38,0.63,0.69}{\textit{#1}}}
\newcommand{\CommentVarTok}[1]{\textcolor[rgb]{0.38,0.63,0.69}{\textbf{\textit{#1}}}}
\newcommand{\ConstantTok}[1]{\textcolor[rgb]{0.53,0.00,0.00}{#1}}
\newcommand{\ControlFlowTok}[1]{\textcolor[rgb]{0.00,0.44,0.13}{\textbf{#1}}}
\newcommand{\DataTypeTok}[1]{\textcolor[rgb]{0.56,0.13,0.00}{#1}}
\newcommand{\DecValTok}[1]{\textcolor[rgb]{0.25,0.63,0.44}{#1}}
\newcommand{\DocumentationTok}[1]{\textcolor[rgb]{0.73,0.13,0.13}{\textit{#1}}}
\newcommand{\ErrorTok}[1]{\textcolor[rgb]{1.00,0.00,0.00}{\textbf{#1}}}
\newcommand{\ExtensionTok}[1]{#1}
\newcommand{\FloatTok}[1]{\textcolor[rgb]{0.25,0.63,0.44}{#1}}
\newcommand{\FunctionTok}[1]{\textcolor[rgb]{0.02,0.16,0.49}{#1}}
\newcommand{\ImportTok}[1]{\textcolor[rgb]{0.00,0.50,0.00}{\textbf{#1}}}
\newcommand{\InformationTok}[1]{\textcolor[rgb]{0.38,0.63,0.69}{\textbf{\textit{#1}}}}
\newcommand{\KeywordTok}[1]{\textcolor[rgb]{0.00,0.44,0.13}{\textbf{#1}}}
\newcommand{\NormalTok}[1]{#1}
\newcommand{\OperatorTok}[1]{\textcolor[rgb]{0.40,0.40,0.40}{#1}}
\newcommand{\OtherTok}[1]{\textcolor[rgb]{0.00,0.44,0.13}{#1}}
\newcommand{\PreprocessorTok}[1]{\textcolor[rgb]{0.74,0.48,0.00}{#1}}
\newcommand{\RegionMarkerTok}[1]{#1}
\newcommand{\SpecialCharTok}[1]{\textcolor[rgb]{0.25,0.44,0.63}{#1}}
\newcommand{\SpecialStringTok}[1]{\textcolor[rgb]{0.73,0.40,0.53}{#1}}
\newcommand{\StringTok}[1]{\textcolor[rgb]{0.25,0.44,0.63}{#1}}
\newcommand{\VariableTok}[1]{\textcolor[rgb]{0.10,0.09,0.49}{#1}}
\newcommand{\VerbatimStringTok}[1]{\textcolor[rgb]{0.25,0.44,0.63}{#1}}
\newcommand{\WarningTok}[1]{\textcolor[rgb]{0.38,0.63,0.69}{\textbf{\textit{#1}}}}
\DefineVerbatimEnvironment{Highlighting}{Verbatim}{commandchars=\\\{\},breaklines,breakanywhere,breaknonspaceingroup,fontsize=\small}
\DefineVerbatimEnvironment{verbatim}{Verbatim}{breaklines,breakanywhere,fontsize=\small}

% 参考资料框
\newtcolorbox{reading}{enhanced, breakable,
  colback=gray!5, colframe=gray!50, boxrule=0.5pt, arc=1.5pt,
  left=7pt, right=7pt, top=5pt, bottom=5pt,
  title={参考资料}, fonttitle=\bfseries\sffamily\small,
  coltitle=black, attach title to upper={\quad}}

% 压轴题框
\newtcolorbox{capstone}[1]{enhanced, breakable,
  colback=accent!4, colframe=accent!70, boxrule=0.7pt, arc=1.5pt,
  left=7pt, right=7pt, top=5pt, bottom=5pt,
  title={#1}, fonttitle=\bfseries\sffamily}

% 回望框
\newtcolorbox{lookback}{enhanced, breakable,
  colback=gray!5, colframe=gray!50, boxrule=0.5pt, arc=1.5pt,
  left=7pt, right=7pt, top=5pt, bottom=5pt,
  title={Editor's Note}, fonttitle=\bfseries\sffamily\small,
  coltitle=black, attach title to upper={\quad},
  before upper app={\setlength{\parskip}{0.5em}}}

% 学员作答
\newtcolorbox{answer}{enhanced, breakable,
  colback=strgreen!3, colframe=strgreen!55, boxrule=0.5pt, arc=1.5pt,
  left=7pt, right=7pt, top=5pt, bottom=5pt,
  title={作答}, fonttitle=\bfseries\sffamily\small,
  coltitle=black, attach title to upper={\quad},
  before upper app={\setlength{\parskip}{0.4em}}}

% 题目标题:\prob{编号}{TYPE};选做 \prob[Optional]{编号}{TYPE};
% 末尾可选参数标注涉及的代码文件
\NewDocumentCommand{\prob}{O{}mmo}{\par\medskip\noindent
  {\bfseries prob #2}\;{\small\sffamily(#3\IfBlankF{#1}{\,·\,#1})}%
  \IfValueT{#4}{\hfill{\footnotesize\ttfamily\color{black!60}#4}}\par\nobreak\smallskip\noindent\ignorespaces}

\begin{document}

\begin{center}
  {\LARGE\bfseries 作业 2:Tensor Core \& Pipeline}\\[6pt]
  {\sffamily Weiming HPC Training Camp \(\times\) LCPU AI Infra Seminars
· Session 3 / 4}
\end{center}

\vspace{0.5em}
\hrule
\vspace{1em}

\setcounter{section}{-1}

\section*{Preface}\label{preface}
\addcontentsline{toc}{section}{Preface}

这次作业主要是 Session 3 的配套练习，内容围绕 Tensor Core 展开，其中
Module 4 还会用到 Session 4 介绍的 tiling、TMA 和
pipeline。完成作业时主要参考课程课件和 PTX ISA 文档。

题型沿用系列惯例，并新增 DERIVE：先手工推导，再写程序验证推导。共七种：
CONCEPT 概念题、DERIVE 推导题、MODIFY 改造题、DEBUG 修 bug 题、
EXPERIMENT 实验题、HANDS-ON 动手题、FROM-SCRATCH 从零实现题。标 Optional
的为选做。涉及代码的题在标题行右侧标出文件路径(相对
\texttt{assignment02/\allowbreak{}})，题面细节以文件头注释为准；FROM-SCRATCH
题。

硬件与构建:主线环境是 B300，\texttt{cuda/\allowbreak{}Makefile} 默认
\texttt{ARCH=\allowbreak{}100f}；M0、M1 也可以在 5090
上完成(\texttt{ARCH=\allowbreak{}120a\ make\ .\allowbreak{}.\allowbreak{}.\allowbreak{}})；
M2 的判测纯 host，无卡可判；M3、M4、M5 需要 B300。本作业统一用显式
\texttt{-\allowbreak{}gencode} 而不用
\texttt{-\allowbreak{}arch=\allowbreak{}sm\_\allowbreak{}XXXa}
简写，原因见
\texttt{assignment02/\allowbreak{}README.\allowbreak{}md}，Makefile
已经配好。

实验纪律:所有性能数字都在自己占满的卡上测；测前测后
\texttt{nvidia-\allowbreak{}smi\ -\allowbreak{}\/\allowbreak{}-\allowbreak{}query-\allowbreak{}compute-\allowbreak{}apps=\allowbreak{}pid，name\ -\allowbreak{}\/\allowbreak{}-\allowbreak{}format=\allowbreak{}csv}
查有没有别的 进程占卡；可能 hang 的程序套
\texttt{timeout\ -\allowbreak{}k\ 5\ \textless{}秒数\textgreater{}}(判测脚本已带)。

AI政策:必做题沿用仓库根目录 \texttt{CLAUDE.\allowbreak{}md}------AI
可以帮你理解、 review，不能替你实现；团队题不设限制。详见
\texttt{assignment02/\allowbreak{}README.\allowbreak{}md}。

\begin{answer}

本次作答环境（2026-09-09）：本地为 AMD Ryzen AI 9 HX 370 / nvcc
12.9.86；本次 \texttt{nvidia-\allowbreak{}smi} 无法连接 NVIDIA
驱动，host/Python 判测在本地完成。GPU 实验使用 SSH
\texttt{b300-\allowbreak{}login} 经 Slurm 分配的 NVIDIA B300 SXM6
\texttt{AC（sm\_\allowbreak{}103，148} SM）/ nvcc 13.0.88，代码使用
\texttt{ARCH=\allowbreak{}100f}。

必做 M0--M6 全部作答；4.4、5.3(d) 和团队 C1/C2
未选做。每题的命令、原始数据和判测输出保存于
\texttt{assignment02/\allowbreak{}results/\allowbreak{}}；路径均相对
\texttt{assignment02/\allowbreak{}}。性能表采用 CUDA events
预热后的平均计时，NCU
数据用于性能归因。\texttt{results/\allowbreak{}README.\allowbreak{}md}
给出证据索引和复现步骤。

\end{answer}

\subsection*{Content}\label{content}
\addcontentsline{toc}{subsection}{Content}

\begin{longtable}[]{@{}
  >{\raggedright\arraybackslash}p{(\linewidth - 6\tabcolsep) * \real{0.2500}}
  >{\raggedright\arraybackslash}p{(\linewidth - 6\tabcolsep) * \real{0.2500}}
  >{\raggedright\arraybackslash}p{(\linewidth - 6\tabcolsep) * \real{0.2500}}
  >{\raggedright\arraybackslash}p{(\linewidth - 6\tabcolsep) * \real{0.2500}}@{}}
\toprule\noalign{}
\begin{minipage}[b]{\linewidth}\raggedright
Module
\end{minipage} & \begin{minipage}[b]{\linewidth}\raggedright
主题
\end{minipage} & \begin{minipage}[b]{\linewidth}\raggedright
对应课件
\end{minipage} & \begin{minipage}[b]{\linewidth}\raggedright
代码位置
\end{minipage} \\
\midrule\noalign{}
\endhead
\bottomrule\noalign{}
\endlastfoot
0 & 环境与峰值 & P1(S008--S021) &
\texttt{cuda/\allowbreak{}m0\_\allowbreak{}env/\allowbreak{}} \\
1 & sm80:fragment 与 mma.sync & P2.2(S025--S041) &
\texttt{cuda/\allowbreak{}m1\_\allowbreak{}sm80/\allowbreak{}} \\
2 & smem 供数:descriptor 与 swizzle & P2.3(S042--S070) &
\texttt{cuda/\allowbreak{}m2\_\allowbreak{}smem/\allowbreak{}} \\
3 & sm100:tcgen05 & P2.4(S071--S093) &
\texttt{cuda/\allowbreak{}m3\_\allowbreak{}tcgen05/\allowbreak{}} \\
4 & 完整 GEMM 四步走 & S084 + Session 4 &
\texttt{cuda/\allowbreak{}m4\_\allowbreak{}gemm/\allowbreak{}} \\
5 & 低精度与 block scaling & P3(S094--S111) &
\texttt{cuda/\allowbreak{}m5\_\allowbreak{}lowprec/\allowbreak{}}、\texttt{kernels/\allowbreak{}} \\
6 & TileLang 对照 & 收尾(S112) & (复用 assignment01) \\
Team & FlashKDA / MSA decode & --- & \texttt{team/\allowbreak{}} \\
\end{longtable}

\section{环境与峰值}\label{ux73afux5883ux4e0eux5cf0ux503c}

先确认工具链和卡都能发出 tensor core 指令：编译一个最小tensor
core程序，测试卡的峰值性能，后面就可以用它来计算各个实现的性能达成率。

\begin{reading}

Session3课件S008-S021；PTX ISA 的 mma 指令一节(形状与 dtype 表)；你手里
每块卡的官方 datasheet(或 whitepaper)。

\end{reading}

\prob{0.1}{HANDS-ON}[cuda/m0\_env/01\_first\_mma.cu]

编译并运行最小 Tensor Core 程序，它使用一个 warp 执行一条 m16n8k16 的
mma.sync 指令，并与 CPU 计算结果进行比较。

\begin{verbatim}
cd assignment02/cuda
make run/m0_env/01_first_mma
\end{verbatim}

在你能使用的 GPU 上分别运行该程序（5090 使用 ARCH=120a make \ldots，B300
使用默认配置）。尝试使用不匹配的 ARCH 编译运行一次，记录现象，并结合
assignment01 Module 8 中 fatbin/JIT 的内容解释原因。可使用 make
\texttt{ptx/\allowbreak{}m0\_\allowbreak{}env/\allowbreak{}01\_\allowbreak{}first\_\allowbreak{}mma}
查看生成的 PTX。

\begin{answer}

硬件：NVIDIA B300 SXM6 AC，compute capability 10.3，148 SM；CUDA 13.0 /
nvcc 13.0.88，\texttt{ARCH=\allowbreak{}100f}。

本次本地 \texttt{nvidia-\allowbreak{}smi}
无法连接驱动（\texttt{results/\allowbreak{}local/\allowbreak{}environment.\allowbreak{}log}），GPU
判测使用 B300，数据来自 Slurm 作业 23618。

\begin{Shaded}
\begin{Highlighting}[]
\BuiltInTok{cd}\NormalTok{ assignment02/cuda}
\FunctionTok{make} \AttributeTok{{-}B}\NormalTok{ run/m0\_env/01\_first\_mma}
\FunctionTok{make} \AttributeTok{{-}B}\NormalTok{ ARCH=89 bin/m0\_env/01\_first\_mma}
\ExtensionTok{./bin/m0\_env/01\_first\_mma}
\end{Highlighting}
\end{Shaded}

匹配架构输出
\texttt{D{[}0{]}{[}0{]}=\allowbreak{}2\ D{[}0{]}{[}7{]}=\allowbreak{}2\ D{[}15{]}{[}0{]}=\allowbreak{}2\ D{[}15{]}{[}7{]}=\allowbreak{}2}、\texttt{PASS}。
不匹配版本的完整错误见
\texttt{results/\allowbreak{}b300/\allowbreak{}arch-\allowbreak{}mismatch.\allowbreak{}log}。Makefile
生成的 fatbin 仅含 \texttt{sm\_\allowbreak{}89} SASS；在
\texttt{sm\_\allowbreak{}103} 上，\texttt{cudaMalloc} 成功，kernel
launch 检查报架构错误。恢复默认架构后强制重编。

PTX/SASS 与 JIT 的区别：SASS 是目标架构的机器码，PTX 可由 driver JIT
编译为目标机器码。driver 从 fatbin 选择镜像，PTX fallback
需要在构建时显式嵌入。

\end{answer}

\prob{0.2}{DERIVE}

推导你所使用 GPU 的 Tensor Core 理论峰值。参考课上 A100
的推导方法（S018--S019），分别计算 5090 和 B300 的 bf16 峰值，并根据
dtype 宽度关系估算 fp8 / fp4 峰值。

开始计算前先明确采用的口径，包括 dense 或 sparse、boost 或 base
频率、FMA 是否计作 2 FLOP 等。完成推导后，再与 datasheet
中的官方数据进行对照；如果结果存在差异，需要说明差异来自哪一项口径。

\begin{longtable}[]{@{}lll@{}}
\toprule\noalign{}
量 & 5090 & B300 \\
\midrule\noalign{}
\endhead
\bottomrule\noalign{}
\endlastfoot
bf16 FLOP/cycle/SM & & \\
bf16 峰值(TFLOPS) & & \\
fp8 峰值(TFLOPS) & & \\
fp4 峰值(TFLOPS) & & \\
datasheet 对照值与口径差异 & & \\
HBM/GDDR 带宽(GB/s) & & \\
机器平衡点(FLOP/byte，bf16) & & \\
\end{longtable}

根据 bf16 峰值和显存带宽计算机器平衡点（FLOP/byte），并与单条 mma
的计算强度（S016，m16n8k16 fp16 为 3.2
FLOP/byte）比较。思考两者之间的差距意味着什么，以及为什么后续 M2--M4
需要从数据供给路径入手优化。

\begin{answer}

硬件：NVIDIA B300 SXM6 AC，compute capability 10.3，148 SM；CUDA 13.0 /
nvcc 13.0.88，\texttt{ARCH=\allowbreak{}100f}。

口径：dense，FMA=2 FLOP，BF16/FP8 均按 FP32 累加。5090 用官方 boost
2.407 GHz；B300 分别按本机
\texttt{clocks.\allowbreak{}max.\allowbreak{}sm=\allowbreak{}2032\ MHz}
计算时钟上界，并列出官方系统规格供对照。

\begin{longtable}[]{@{}
  >{\raggedright\arraybackslash}p{(\linewidth - 4\tabcolsep) * \real{0.2727}}
  >{\raggedleft\arraybackslash}p{(\linewidth - 4\tabcolsep) * \real{0.3636}}
  >{\raggedleft\arraybackslash}p{(\linewidth - 4\tabcolsep) * \real{0.3636}}@{}}
\toprule\noalign{}
\begin{minipage}[b]{\linewidth}\raggedright
量
\end{minipage} & \begin{minipage}[b]{\linewidth}\raggedleft
RTX 5090（资料对照）
\end{minipage} & \begin{minipage}[b]{\linewidth}\raggedleft
本次 B300
\end{minipage} \\
\midrule\noalign{}
\endhead
\bottomrule\noalign{}
\endlastfoot
SM 数 & 170 & 实测 148 \\
BF16 FLOP/cycle/SM & 512 & 8192 \\
BF16 按上述时钟算 TFLOPS & 209.5 & 2463.6 \\
FP8 纯位宽估计 TFLOPS & 419.0 & 4927.3 \\
FP4 纯位宽估计 TFLOPS & 838.0 & 9854.6 \\
官方 dense BF16 / FP8 / FP4 TFLOPS & 209.5 / 419 / 1676 & 2250 / 4500 /
13500 \\
显存带宽 GB/s & 1792 & 8000 \\
官方 BF16 机器平衡点 FLOP/B & 116.9 & 281.25 \\
最大时钟上界对应平衡点 FLOP/B & 116.9 & 307.95 \\
\end{longtable}

推导：\(P=N_{\mathrm{SM}}qf\)。5090 为
\(170\times512\times2.407=209500.16\) GFLOPS；B300 的时钟上界估算为
\(148\times8192\times2.032=2463629.312\) GFLOPS。由官方 2250 TFLOPS
反推，等效重载频率约为 1.856 GHz。

官方 HGX B300 的 BF16/FP8 是八卡 sparse 数据：36/72 PFLOPS 除以 8 再除以
2，得到 2250/4500 TFLOPS。FP4 写作 sparse\textbar dense=144\textbar108
PFLOPS，dense 单卡为 13500。5090 的 FP4 相对 BF16 FP32 累加为 8 倍；B300
Ultra FP4 也有额外吞吐增强，因此 FP4 峰值采用官方规格。

后续 4.5 按官方 2250 TFLOPS、8000 GB/s 统一计算 roofline。单条 MMA 的
3.2 FLOP/B
远低于两块卡的平衡点；需要跨指令复用、减少搬运和隐藏供数延迟。

来源：\href{https://images.nvidia.com/aem-dam/Solutions/geforce/blackwell/nvidia-rtx-blackwell-gpu-architecture.pdf}{RTX
Blackwell 白皮书 Table
3}、\href{https://www.nvidia.com/en-us/data-center/hgx/}{HGX 规格及
dense/sparse
脚注}、\href{https://gfxcourses.stanford.edu/cs149/fall25content/media/proghardware/11_SpecializedHardwareProgramming.pdf}{Stanford
CS149 Tensor Core 吞吐表}。本机配置见
\texttt{results/\allowbreak{}b300/\allowbreak{}device.\allowbreak{}log}、\texttt{gpu-\allowbreak{}details.\allowbreak{}csv}。

\end{answer}

\prob{0.3}{CONCEPT}

判断下列说法是否正确，并给出一句理由。

\begin{enumerate}
\def\labelenumi{(\alph{enumi})}
\item
  一条 mma 的计算强度，分子是 \(2MNK\)，分母按 A、B 读入与 D 写回
  的字节总和计(S016 的口径)。
\item
  mma.sync 是 warp 级协作指令:32 个 lane 各持 fragment 的一部分， 要求全
  warp 一致地执行这条指令；有 lane 发散时行为未定义。
\item
  增大 mma 的形状 M/N/K 能提高单条指令的计算强度，而且没有代价，
  所以指令形状越大越好。
\item
  只要单条 mma 的计算强度低于机器平衡点，GEMM kernel 就不可能逼近
  计算峰值。
\end{enumerate}

\begin{answer}

硬件：NVIDIA B300 SXM6 AC，compute capability 10.3，148 SM；CUDA 13.0 /
nvcc 13.0.88，\texttt{ARCH=\allowbreak{}100f}。

\begin{enumerate}
\def\labelenumi{(\alph{enumi})}
\item
  对，按题设口径分母包括 A/B 读取及 D 写回，分子为 \(2MNK\)。
\item
  对，\texttt{mma.\allowbreak{}sync.\allowbreak{}aligned} 要求参与 warp
  的 lane 一致执行，fragment 是各 lane 合作持有的矩阵切片。
\item
  错，更大的形状增加寄存器或 TMEM、shared memory
  和同步开销，还可能增加边界浪费，降低可驻留 CTA 数。
\item
  错，kernel 可在 shared memory 中复用 A/B，并让累加结果跨多条 MMA
  留在寄存器或 TMEM。沿 K 保留累加器减少中间输出流量，沿 M/N
  复用输入减少重复读取，kernel 的强度可以远高于单条指令。
\end{enumerate}

\end{answer}

\section{sm80:fragment 与 mma.sync}\label{sm80fragment-ux4e0e-mma.sync}

课上对 m16n8k16 fp16 推过 fragment 公式、手搓过单 tile mma
(C03--C07)。现在你手推一遍这个 m16n8k32 fp8 shape， 再看 ldmatrix
到底发挥了什么作用。

\begin{reading}

课件 S025--S041、C03--C08；PTX ISA 的 ``Matrix Fragments for
mma.m16n8k32'' 一节与 ``Warp-level matrix load instruction: ldmatrix''
一节；fp8 转换用
\texttt{cuda\_\allowbreak{}fp8.\allowbreak{}h}(\texttt{\_\allowbreak{}\_\allowbreak{}nv\_\allowbreak{}fp8\_\allowbreak{}e4m3})。

\end{reading}

\prob{1.1}{DERIVE}[cuda/m1\_sm80/01\_fragment\_map.cu]

根据 PTX 文档中的 fragment 布局，推导
\texttt{mma.\allowbreak{}sync.\allowbreak{}aligned.\allowbreak{}m16n8k32.\allowbreak{}row.\allowbreak{}col.\allowbreak{}f32.\allowbreak{}e4m3.\allowbreak{}e4m3.\allowbreak{}f32}
对应的 A、B fragment 映射公式，并完成文件中的四个函数。 判测使用纯 host
的 32-lane 真值表比对，无需 GPU 即可完成：

\begin{verbatim}
cd assignment02/cuda
make run/m1_sm80/01_fragment_map
\end{verbatim}

\begin{answer}

硬件：本地 AMD Ryzen AI 9 HX 370（CPU 实验，compute capability
不适用）；nvcc 12.9.86。

设 \texttt{lane} 为 lane 号，\texttt{i} 为本 lane fragment 中的 byte
序号， \texttt{r\ =\allowbreak{}\ i\ /\allowbreak{}\ 4} 为 b32
寄存器号，\texttt{j\ =\allowbreak{}\ i\ \%\ 4} 为寄存器内 byte 序号。

对于
A（\texttt{i=\allowbreak{}0.\allowbreak{}.\allowbreak{}15}，逻辑形状
\(16\times32\)）：

\[
\begin{aligned}
\mathrm{row}_A(lane,i) &= (lane\mathbin{>>}2)+8(r\mathbin{\&}1),\\
\mathrm{col}_A(lane,i) &= 4(lane\mathbin{\&}3)+16(r\mathbin{>>}1)+j.
\end{aligned}
\]

\begin{Shaded}
\begin{Highlighting}[]
\AttributeTok{static} \DataTypeTok{int}\NormalTok{ a\_row\_of}\OperatorTok{(}\DataTypeTok{int}\NormalTok{ lane}\OperatorTok{,} \DataTypeTok{int}\NormalTok{ i}\OperatorTok{)} \OperatorTok{\{}
    \DataTypeTok{int}\NormalTok{ r }\OperatorTok{=}\NormalTok{ i }\OperatorTok{/} \DecValTok{4}\OperatorTok{;}
    \ControlFlowTok{return} \OperatorTok{(}\NormalTok{lane }\OperatorTok{\textgreater{}\textgreater{}} \DecValTok{2}\OperatorTok{)} \OperatorTok{+} \DecValTok{8} \OperatorTok{*} \OperatorTok{(}\NormalTok{r }\OperatorTok{\&} \DecValTok{1}\OperatorTok{);}
\OperatorTok{\}}

\AttributeTok{static} \DataTypeTok{int}\NormalTok{ a\_col\_of}\OperatorTok{(}\DataTypeTok{int}\NormalTok{ lane}\OperatorTok{,} \DataTypeTok{int}\NormalTok{ i}\OperatorTok{)} \OperatorTok{\{}
    \DataTypeTok{int}\NormalTok{ r }\OperatorTok{=}\NormalTok{ i }\OperatorTok{/} \DecValTok{4}\OperatorTok{;}
    \ControlFlowTok{return} \DecValTok{4} \OperatorTok{*} \OperatorTok{(}\NormalTok{lane }\OperatorTok{\&} \DecValTok{3}\OperatorTok{)} \OperatorTok{+} \DecValTok{16} \OperatorTok{*} \OperatorTok{(}\NormalTok{r }\OperatorTok{\textgreater{}\textgreater{}} \DecValTok{1}\OperatorTok{)} \OperatorTok{+} \OperatorTok{(}\NormalTok{i }\OperatorTok{\&} \DecValTok{3}\OperatorTok{);}
\OperatorTok{\}}
\end{Highlighting}
\end{Shaded}

对于 B（\texttt{i=\allowbreak{}0.\allowbreak{}.\allowbreak{}7}，逻辑形状
\(32\times8\)，函数返回值的 row 就是 K 维）：

\[
\begin{aligned}
\mathrm{row}_B(lane,i) &= 4(lane\mathbin{\&}3)+(i\mathbin{\&}3)+16(i\mathbin{>>}2),\\
\mathrm{col}_B(lane,i) &= lane\mathbin{>>}2.
\end{aligned}
\]

\begin{Shaded}
\begin{Highlighting}[]
\AttributeTok{static} \DataTypeTok{int}\NormalTok{ b\_row\_of}\OperatorTok{(}\DataTypeTok{int}\NormalTok{ lane}\OperatorTok{,} \DataTypeTok{int}\NormalTok{ i}\OperatorTok{)} \OperatorTok{\{}
    \ControlFlowTok{return} \DecValTok{4} \OperatorTok{*} \OperatorTok{(}\NormalTok{lane }\OperatorTok{\&} \DecValTok{3}\OperatorTok{)} \OperatorTok{+} \OperatorTok{(}\NormalTok{i }\OperatorTok{\&} \DecValTok{3}\OperatorTok{)} \OperatorTok{+} \DecValTok{16} \OperatorTok{*} \OperatorTok{(}\NormalTok{i }\OperatorTok{\textgreater{}\textgreater{}} \DecValTok{2}\OperatorTok{);}
\OperatorTok{\}}

\AttributeTok{static} \DataTypeTok{int}\NormalTok{ b\_col\_of}\OperatorTok{(}\DataTypeTok{int}\NormalTok{ lane}\OperatorTok{,} \DataTypeTok{int}\NormalTok{ i}\OperatorTok{)} \OperatorTok{\{}
    \ControlFlowTok{return}\NormalTok{ lane }\OperatorTok{\textgreater{}\textgreater{}} \DecValTok{2}\OperatorTok{;}
\OperatorTok{\}}
\end{Highlighting}
\end{Shaded}

例如 lane 0 的 A fragment 是
\texttt{(row,\allowbreak{}col)=\allowbreak{}(0,\allowbreak{}0.\allowbreak{}.\allowbreak{}3),\allowbreak{}(8,\allowbreak{}0.\allowbreak{}.\allowbreak{}3),\allowbreak{}(0,\allowbreak{}16.\allowbreak{}.\allowbreak{}19),\allowbreak{}(8,\allowbreak{}16.\allowbreak{}.\allowbreak{}19)}；
lane 0 的 B fragment 是
\texttt{(k,\allowbreak{}n)=\allowbreak{}(0.\allowbreak{}.\allowbreak{}3,\allowbreak{}0),\allowbreak{}(16.\allowbreak{}.\allowbreak{}19,\allowbreak{}0)}。本次运行
\texttt{ARCH=\allowbreak{}89\ make\ -\allowbreak{}B\ run/\allowbreak{}m1\_\allowbreak{}sm80/\allowbreak{}01\_\allowbreak{}fragment\_\allowbreak{}map}
输出 \texttt{PASS}，记录见
\texttt{results/\allowbreak{}local/\allowbreak{}host-\allowbreak{}tests.\allowbreak{}log}。

附加问题：A 的同一个 b32 寄存器含有 4 个连续 FP8 byte，公式中变化的
\texttt{j} 对应 A 的列方向，即 K 方向。\texttt{ldmatrix}
按矩阵行地址组织 shared memory 访问；各 lane 按 PTX fragment
映射提供地址，将相应的连续 FP8 元素装入寄存器。

\end{answer}

\prob{1.2}{DEBUG}[cuda/m1\_sm80/02\_bug\_fragment.cu]

这个程序发一条 m16n8k16 fp16 mma，判测会 FAIL。先运行一遍，后改动:

\begin{enumerate}
\def\labelenumi{(\alph{enumi})}
\item
  描述症状：D 的哪些位置错、错成了什么(和对的部分是什么关系)；
\item
  修好它，并解释错的是哪个 fragment
  的哪部分映射，为什么恰好产生(a)的症状。
\end{enumerate}

\begin{verbatim}
cd assignment02/cuda
make run/m1_sm80/02_bug_fragment
\end{verbatim}

\begin{answer}

硬件：NVIDIA B300 SXM6 AC，compute capability 10.3，148 SM；CUDA 13.0 /
nvcc 13.0.88，\texttt{ARCH=\allowbreak{}100f}。

先运行未修改版本：前四处错误是
\texttt{D{[}8{]}{[}0.\allowbreak{}.\allowbreak{}3{]}}，got 为
\texttt{-\allowbreak{}1,\allowbreak{}-\allowbreak{}5,\allowbreak{}-\allowbreak{}16,\allowbreak{}15}，want
为
\texttt{-\allowbreak{}11,\allowbreak{}5,\allowbreak{}-\allowbreak{}14,\allowbreak{}9}；总计
\texttt{59\ /\allowbreak{}\ 128\ mismatches}。上半 8 行正确，下半 8
行重复上半的结果，其中 5 个元素碰巧相等。

A 的 \texttt{a2/\allowbreak{}a3/\allowbreak{}a6/\allowbreak{}a7} 应来自
\texttt{group+8} 行，错误版本仍读取 \texttt{group} 行；B 和 D
的映射无误。修复这四个源地址后，\texttt{make\ -\allowbreak{}B\ run/\allowbreak{}m1\_\allowbreak{}sm80/\allowbreak{}02\_\allowbreak{}bug\_\allowbreak{}fragment}
输出 \texttt{PASS}。原始源码在
\texttt{results/\allowbreak{}original/\allowbreak{}02\_\allowbreak{}bug\_\allowbreak{}fragment.\allowbreak{}cu}，前后记录在
\texttt{results/\allowbreak{}b300/\allowbreak{}debug-\allowbreak{}before.\allowbreak{}log}
和
\texttt{02\_\allowbreak{}bug\_\allowbreak{}fragment.\allowbreak{}log}。

\end{answer}

\begin{capstone}{prob 1.3(FROM-SCRATCH):手写单 tile fp8 mma}

在
\texttt{cuda/\allowbreak{}m1\_\allowbreak{}sm80/\allowbreak{}03\_\allowbreak{}mma\_\allowbreak{}fp8.\allowbreak{}cu}
中从零实现一个单 tile 的 fp8 mma，不提供代码骨架。使用 \texttt{m16n8k32}
e4m3 mma、f32 累加，并手动完成 fragment 装载（本题不使用
ldmatrix），最后与 CPU 参考结果进行严格相等比较。

要求如下：

输入使用小整数，保证转换为 e4m3 后可以精确表示；

程序接受一个 seed 参数，例如
\texttt{.\allowbreak{}/\allowbreak{}prog\ 123}；

输出必须以 \texttt{PASS} 或 \texttt{MISMATCH} / \texttt{FAIL} 开头；

fp8 与 float 的互转直接使用
\texttt{cuda\_\allowbreak{}fp8.\allowbreak{}h}。

1.1 中推导的 fragment
映射公式会直接用在这里。映射只要有一处错误，随机数据的判测就无法通过。

判测脚本会使用五个 seed，全部通过才算完成：

\begin{verbatim}
cd assignment02/cuda/m1_sm80
./judge_mma_fp8.sh 03_mma_fp8.cu
\end{verbatim}

\end{capstone}

\begin{answer}

硬件：NVIDIA B300 SXM6 AC，compute capability 10.3，148 SM；CUDA 13.0 /
nvcc 13.0.88，\texttt{ARCH=\allowbreak{}100f}。

实现文件
\texttt{cuda/\allowbreak{}m1\_\allowbreak{}sm80/\allowbreak{}03\_\allowbreak{}mma\_\allowbreak{}fp8.\allowbreak{}cu}。按
1.1 的映射从 global memory 逐 byte 装载，四个 A 寄存器、两个 B
寄存器组成一条 \texttt{m16n8k32} E4M3 MMA；D 按 m16n8
的四元素映射写回。随机小整数都可被 FP8 精确表示，f32
整数和也可精确表示，因而使用严格相等比较。

\begin{Shaded}
\begin{Highlighting}[]
\BuiltInTok{cd}\NormalTok{ assignment02/cuda/m1\_sm80}
\FunctionTok{bash}\NormalTok{ judge\_mma\_fp8.sh 03\_mma\_fp8.cu}
\end{Highlighting}
\end{Shaded}

seed
\texttt{1,\allowbreak{}7,\allowbreak{}42,\allowbreak{}1234,\allowbreak{}99999}
全部
\texttt{PASS\ .\allowbreak{}.\allowbreak{}.\allowbreak{}\ bad=\allowbreak{}0}，最终
\texttt{JUDGE:\allowbreak{}\ PASS}；原始输出
\texttt{results/\allowbreak{}b300/\allowbreak{}judge-\allowbreak{}fp8.\allowbreak{}log}。

\end{answer}

\prob{1.4}{MODIFY}[cuda/m1\_sm80/04\_ldmatrix.cu]

在给定骨架中保留 1.3 的手工装载方式，并另外实现一条使用
\texttt{ldmatrix} 的装载路径。两种实现需要共存，并分别通过判测。

根据 PTX 文档选择合适的 \texttt{ldmatrix}
变体（\texttt{.\allowbreak{}x1/\allowbreak{}.\allowbreak{}x2/\allowbreak{}.\allowbreak{}x4}，以及是否使用
\texttt{.\allowbreak{}trans}）。B 矩阵在 shared memory
中的布局还需要满足 \texttt{ldmatrix} 的 16 byte
行地址要求，具体约束见文件头说明。

两条路径都通过判测后，使用
\texttt{make\ ptx/\allowbreak{}m1\_\allowbreak{}sm80/\allowbreak{}04\_\allowbreak{}ldmatrix}
或 \texttt{nvdisasm} 查看生成的指令，分别统计 \texttt{smem\ →\ fragment}
阶段的装载指令和地址计算指令数量，并回答：

\begin{enumerate}
\def\labelenumi{(\alph{enumi})}
\item
  \texttt{ldmatrix} 省掉了手工装载中的哪些工作？
\item
  为什么这些工作在手工装载路径中无法避免？
\end{enumerate}

\begin{answer}

硬件：NVIDIA B300 SXM6 AC，compute capability 10.3，148 SM；CUDA 13.0 /
nvcc 13.0.88，\texttt{ARCH=\allowbreak{}100f}。

\texttt{make\ -\allowbreak{}B\ run/\allowbreak{}m1\_\allowbreak{}sm80/\allowbreak{}04\_\allowbreak{}ldmatrix}：seeds=1、7、42
的 manual/ldsm 两条路径均 \texttt{PASS(0)}。A 使用
\texttt{ldmatrix.\allowbreak{}m8n8.\allowbreak{}x4.\allowbreak{}b16}，B
使用 \texttt{.\allowbreak{}x2.\allowbreak{}b16}，均不转置；B 改用
\texttt{{[}N{]}{[}K{]}} shared 布局，让每行的 16B 连续 FP8
满足行地址要求。b16 只是 16 个原始 bit，实际装入两个连续 E4M3。

\begin{Shaded}
\begin{Highlighting}[]
\FunctionTok{make} \AttributeTok{{-}B}\NormalTok{ ptx/m1\_sm80/04\_ldmatrix}
\ExtensionTok{cuobjdump} \AttributeTok{{-}{-}dump{-}sass}\NormalTok{ bin/m1\_sm80/04\_ldmatrix}
\end{Highlighting}
\end{Shaded}

计数范围限定在生成 PTX 的
\texttt{bar.\allowbreak{}warp.\allowbreak{}sync}
之后、\texttt{mma.\allowbreak{}sync} 之前，排除 global
staging、输出地址和常量初始化。

\begin{longtable}[]{@{}
  >{\raggedright\arraybackslash}p{(\linewidth - 8\tabcolsep) * \real{0.1579}}
  >{\raggedleft\arraybackslash}p{(\linewidth - 8\tabcolsep) * \real{0.2105}}
  >{\raggedleft\arraybackslash}p{(\linewidth - 8\tabcolsep) * \real{0.2105}}
  >{\raggedleft\arraybackslash}p{(\linewidth - 8\tabcolsep) * \real{0.2105}}
  >{\raggedleft\arraybackslash}p{(\linewidth - 8\tabcolsep) * \real{0.2105}}@{}}
\toprule\noalign{}
\begin{minipage}[b]{\linewidth}\raggedright
路径
\end{minipage} & \begin{minipage}[b]{\linewidth}\raggedleft
shared 装载指令
\end{minipage} & \begin{minipage}[b]{\linewidth}\raggedleft
地址算术/位运算
\end{minipage} & \begin{minipage}[b]{\linewidth}\raggedleft
shared 符号地址 mov
\end{minipage} & \begin{minipage}[b]{\linewidth}\raggedleft
byte 打包算术
\end{minipage} \\
\midrule\noalign{}
\endhead
\bottomrule\noalign{}
\endlastfoot
manual & 4×u32 + 8×u8 = 12 & 12 & 2 & 12 \\
ldmatrix & x4 + x2 = 2 & 10 & 2 & 0 \\
\end{longtable}

源码逐 byte 描述 A，编译器将相邻的四个 byte 合并为 u32；B 的 K-major
手工路径跨 8B 取字节，需要装载、移位与 OR 打包。ldmatrix 通过 warp
协作完成分发和打包，由线程提供矩阵行地址。手工路径也可通过预转置 B
减少普通 load。

证据：\texttt{results/\allowbreak{}b300/\allowbreak{}04\_\allowbreak{}ldmatrix.\allowbreak{}log}、\texttt{ldmatrix.\allowbreak{}ptx.\allowbreak{}txt}、\texttt{ldmatrix.\allowbreak{}sass.\allowbreak{}txt}、\texttt{ldmatrix-\allowbreak{}counts.\allowbreak{}txt}。

\end{answer}

\prob{1.5}{EXPERIMENT}[cuda/m1\_sm80/05\_ldsm\_stride.cu]

观察行跨度对 \texttt{ldmatrix} 的影响。对同一个 16×16 fp16
tile，分别使用 32 B、64 B、128 B 和 128+16 B（padding）四种行跨度。

先根据课上介绍的 bank 模型，预测四种情况下执行一次 \texttt{ldmatrix}
所需的 wavefront 数及其比例，再运行程序并使用 Nsight Compute 测量：

\begin{Shaded}
\begin{Highlighting}[]
\BuiltInTok{cd}\NormalTok{ assignment02/cuda}
\FunctionTok{make}\NormalTok{ run/m1\_sm80/05\_ldsm\_stride}
\ExtensionTok{ncu} \AttributeTok{{-}{-}metrics}\NormalTok{ l1tex\_\_data\_pipe\_lsu\_wavefronts\_mem\_shared\_op\_ld.sum,l1tex\_\_data\_bank\_conflicts\_pipe\_lsu\_mem\_shared\_op\_ld.sum ./bin/m1\_sm80/05\_ldsm\_stride}
\end{Highlighting}
\end{Shaded}

\begin{longtable}[]{@{}
  >{\raggedright\arraybackslash}p{(\linewidth - 8\tabcolsep) * \real{0.2000}}
  >{\raggedright\arraybackslash}p{(\linewidth - 8\tabcolsep) * \real{0.2000}}
  >{\raggedright\arraybackslash}p{(\linewidth - 8\tabcolsep) * \real{0.2000}}
  >{\raggedright\arraybackslash}p{(\linewidth - 8\tabcolsep) * \real{0.2000}}
  >{\raggedright\arraybackslash}p{(\linewidth - 8\tabcolsep) * \real{0.2000}}@{}}
\toprule\noalign{}
\begin{minipage}[b]{\linewidth}\raggedright
档位
\end{minipage} & \begin{minipage}[b]{\linewidth}\raggedright
预测 wavefront 比
\end{minipage} & \begin{minipage}[b]{\linewidth}\raggedright
实测 wavefront
\end{minipage} & \begin{minipage}[b]{\linewidth}\raggedright
实测 conflict
\end{minipage} & \begin{minipage}[b]{\linewidth}\raggedright
平均 cycle
\end{minipage} \\
\midrule\noalign{}
\endhead
\bottomrule\noalign{}
\endlastfoot
32 B & & & & \\
64 B & & & & \\
128 B & & & & \\
128+16 B & & & & \\
\end{longtable}

比较预测与实测结果：哪一种行跨度使 wavefront 数增加到 4 倍？ wavefront
的比例应与 bank 模型一致，但实际耗时的差距通常没有这么大。 结合 8 个
warp 的占用情况，解释为什么 wavefront 增加 4 倍并不会使总耗时也增加 4
倍。

\begin{answer}

硬件：NVIDIA B300 SXM6 AC，compute capability 10.3，148 SM；CUDA 13.0 /
nvcc 13.0.88，\texttt{ARCH=\allowbreak{}100f}。

运行题面
\texttt{make\ -\allowbreak{}B\ run/\allowbreak{}m1\_\allowbreak{}sm80/\allowbreak{}05\_\allowbreak{}ldsm\_\allowbreak{}stride}
与两个 NCU metrics，计数采用每档第二次 launch（与第一次相同），未经
profiler 的 cycle 数据用于时延比较。

\begin{longtable}[]{@{}
  >{\raggedright\arraybackslash}p{(\linewidth - 8\tabcolsep) * \real{0.1579}}
  >{\raggedleft\arraybackslash}p{(\linewidth - 8\tabcolsep) * \real{0.2105}}
  >{\raggedleft\arraybackslash}p{(\linewidth - 8\tabcolsep) * \real{0.2105}}
  >{\raggedleft\arraybackslash}p{(\linewidth - 8\tabcolsep) * \real{0.2105}}
  >{\raggedleft\arraybackslash}p{(\linewidth - 8\tabcolsep) * \real{0.2105}}@{}}
\toprule\noalign{}
\begin{minipage}[b]{\linewidth}\raggedright
行跨度
\end{minipage} & \begin{minipage}[b]{\linewidth}\raggedleft
预测 wavefront 比（32B=1）
\end{minipage} & \begin{minipage}[b]{\linewidth}\raggedleft
实测 wavefront
\end{minipage} & \begin{minipage}[b]{\linewidth}\raggedleft
实测 conflict
\end{minipage} & \begin{minipage}[b]{\linewidth}\raggedleft
平均 cycle
\end{minipage} \\
\midrule\noalign{}
\endhead
\bottomrule\noalign{}
\endlastfoot
32 B & 1 & 16384 & 8192 & 9.72 \\
64 B & 2 & 32768 & 24576 & 10.75 \\
128 B & 4 & 65536 & 57344 & 16.06 \\
144 B & 0.5 & 8192 & 0 & 9.23 \\
\end{longtable}

按 32 个 4B bank，行地址步长分别为 8、16、32、36 bank；一组八行的起始
16B chunk 分别覆盖 4、2、1、8 个不同位置，因而相对无冲突 144B
档的冲突倍数为 2:4:8:1。128B 是 32B 的 4 倍，但 cycle 仅约 1.65 倍。

wavefront 统计 shared 服务次数。程序发射 8 个 warp，调度可重叠各 warp
的地址计算、循环、XOR 与访存等待，kernel
总耗时由这些开销共同决定。表中计数为整个 kernel 的总量，原始记录见
\texttt{05\_\allowbreak{}ldsm\_\allowbreak{}stride.\allowbreak{}log}、\texttt{ncu-\allowbreak{}stride.\allowbreak{}csv}。

\end{answer}

\begin{lookback}

1.3--1.5 关注的是同一个问题：怎样把数据从 shared memory 装入 Tensor Core
的 fragment。1.3 手动完成每个 lane 的装载，1.4 使用 \texttt{ldmatrix}
简化这一步，而 1.5 说明即使使用 \texttt{ldmatrix}，shared memory 的 bank
conflict 仍然会影响实际效率。

后面的模块会继续沿着这条数据供给路径展开：M2 中使用 descriptor
描述布局，M4 中进一步使用 TMA 和 pipeline 完成数据搬运。

\end{lookback}

\section{smem 供数:descriptor 与
swizzle}\label{smem-ux4f9bux6570descriptor-ux4e0e-swizzle}

sm90 起，tensor core 从 smem 取数不再经过 fragment 装载指令，而是 读一个
64 位 descriptor 按 canonical 布局自取。descriptor 与 swizzle 的格式在
sm100 的 tcgen05 上原样沿用，本模块推的每个公式都是 M3、 M4 的直接组件。

\begin{reading}

课件 S042--S070(proxy:S051；core matrix 与 LBO/SBO:S057--S061；
swizzle:S062--S065)；PTX ISA 的 ``Asynchronous Warpgroup MMA Shared
Memory Layout'' 与 swizzling 小节。

\end{reading}

\prob{2.1}{CONCEPT}

\begin{enumerate}
\def\labelenumi{(\alph{enumi})}
\tightlist
\item
  一个 warpgroup 使用 wgmma 读取刚写入 shared memory
  的数据。将下面六个操作排成正确顺序，并说明每一步用于避免哪两个参与者之间的哪种乱序：
\end{enumerate}

\texttt{wgmma.\allowbreak{}mma\_\allowbreak{}async} /
\texttt{st.\allowbreak{}shared} /
\texttt{wgmma.\allowbreak{}commit\_\allowbreak{}group} /
\texttt{fence.\allowbreak{}proxy.\allowbreak{}async} /
\texttt{wgmma.\allowbreak{}fence} /
\texttt{wgmma.\allowbreak{}wait\_\allowbreak{}group}

\begin{enumerate}
\def\labelenumi{(\alph{enumi})}
\setcounter{enumi}{1}
\tightlist
\item
  判断下列说法是否正确，并给出一句理由。
\end{enumerate}

\begin{enumerate}
\def\labelenumi{\arabic{enumi}.}
\tightlist
\item
  \texttt{fence.\allowbreak{}proxy.\allowbreak{}async} 是 wgmma
  专属的指令，TMA 与 tcgen05 的场景不需要它。
\item
  \texttt{wgmma.\allowbreak{}commit\_\allowbreak{}group}
  会阻塞，直到它之前发射的 wgmma 全部完成。
\item
  不加 \texttt{fence.\allowbreak{}proxy.\allowbreak{}async} 时，wgmma
  可能读到 shared memory 中的旧值，因为 \texttt{st.\allowbreak{}shared}
  的写经过 generic proxy，而 wgmma 的读经过 async proxy。
\end{enumerate}

\begin{answer}

硬件：本地 AMD Ryzen AI 9 HX 370（CPU 实验，compute capability
不适用）；nvcc 12.9.86。

\begin{enumerate}
\def\labelenumi{(\alph{enumi})}
\tightlist
\item
  顺序为：\texttt{st.\allowbreak{}shared\ →\ fence.\allowbreak{}proxy.\allowbreak{}async.\allowbreak{}shared:\allowbreak{}:\allowbreak{}cta\ →\ wgmma.\allowbreak{}fence\ →\ wgmma.\allowbreak{}mma\_\allowbreak{}async\ →\ wgmma.\allowbreak{}commit\_\allowbreak{}group\ →\ wgmma.\allowbreak{}wait\_\allowbreak{}group\ 0}。若
  shared 数据由多个线程提供，还必须在消费者使用前用适当的 CTA/warpgroup
  同步确保所有生产者已完成。
\end{enumerate}

\texttt{st.\allowbreak{}shared} 产生 generic-proxy 写；proxy fence 使
async-proxy 的 MMA 能观察它；\texttt{wgmma.\allowbreak{}fence}
将此前的寄存器访问与后续 WGMMA 排序；MMA 异步发射，commit
将其组成待完成组，wait 等待相应结果就绪。

\begin{enumerate}
\def\labelenumi{(\alph{enumi})}
\setcounter{enumi}{1}
\tightlist
\item
  1 错：跨 generic/async proxy 的 TMA、tcgen05
  场景同样需要按数据依赖使用 proxy fence。2 错：commit
  完成分组后返回，结果就绪由 wait 保证。3 对：缺失跨 proxy
  排序时，异步消费者可能读取 shared 中的旧值。
\end{enumerate}

出处：PTX ISA 的 WGMMA asynchronous operations / proxy fence 小节。

\end{answer}

\prob{2.2}{DERIVE}[cuda/m2\_smem/02\_descriptor.cu]

根据文件头给出的位域，实现 SM100 的 64 位 smem matrix descriptor
编码函数，并分别对下面三种情况推导 LBO、SBO 和 layout：

\begin{itemize}
\tightlist
\item
  K-major，无 swizzle；
\item
  K-major，128B swizzle；
\item
  MN-major，128B swizzle。
\end{itemize}

判测为纯 host，无需 GPU。测试使用的描述符真值已经在 B300 上通过实际
tcgen05 指令验证，3.2 中也会使用这三组描述符：

\begin{verbatim}
cd assignment02/cuda
make run/m2_smem/02_descriptor
\end{verbatim}

场景 2 和场景 3 最终得到的 descriptor 相同。在报告中回答：MN-major 与
K-major 的区别体现在哪里？

\begin{answer}

硬件：本地 AMD Ryzen AI 9 HX 370（CPU 实验，compute capability
不适用）；nvcc 12.9.86。

编码公式（所有地址以 byte 输入）：

\[\begin{aligned}
d={}&((a\gg4)\mathbin{\&}16383)\;|\;(((L\gg4)\mathbin{\&}16383)\ll16)\\
 &{}|\;(((S\gg4)\mathbin{\&}16383)\ll32)\;|\;(1\ll46)\;|\;(layout\ll61).
\end{aligned}\]

\begin{longtable}[]{@{}lrrrl@{}}
\toprule\noalign{}
场景 & LBO (B) & SBO (B) & layout & descriptor \\
\midrule\noalign{}
\endhead
\bottomrule\noalign{}
\endlastfoot
K-major NONE & 128 & 1024 & 0 & \texttt{0x0000404000080100} \\
K-major 128B & 0 & 1024 & 2 & \texttt{0x4000404000000200} \\
MN-major 128B & 0 & 1024 & 2 & \texttt{0x4000404000000300} \\
\end{longtable}

无 swizzle 的 core matrix 为 \(8\times16\) B，沿 K 移一个 core 为 128
B；64 个 bf16 的 K 维有 8 个 core，沿 MN 跨 core 为 1024 B。128B swizzle
的 atom 为 \(8\times128=1024\) B，LBO 被忽略，填 0。

场景 2/3 的布局字段相同，起始地址字段不同，因此完整的 64
位描述符值不同。K-major/MN-major 的转置解释在 MMA 的 instruction
descriptor 中，数据按对应 major 方向摆放。

\texttt{ARCH=\allowbreak{}89\ make\ -\allowbreak{}B\ run/\allowbreak{}m2\_\allowbreak{}smem/\allowbreak{}02\_\allowbreak{}descriptor}：三个场景均
\texttt{PASS}；见
\texttt{results/\allowbreak{}local/\allowbreak{}host-\allowbreak{}tests.\allowbreak{}log}。

\end{answer}

\prob{2.3}{FROM-SCRATCH}[cuda/m2\_smem/03\_swizzle.cu]

实现 128B、64B 和 32B 三种 swizzle 模式的地址映射函数，将逻辑坐标映射到
atom 内的物理字节偏移。

课件 S064 已经推导了 128B swizzle 的地址位异或关系；64B 和 32B
两种模式需要根据 PTX ISA 的 swizzling 一节自行推导。

判测分为两部分：首先检查映射是否为双射，然后检查列访问是否存在 bank
conflict。单纯使用 padding
或恒等映射虽然可能通过第一项，但无法通过第二项。判测同样为纯 host：

\begin{verbatim}
cd assignment02/cuda
make run/m2_smem/03_swizzle
\end{verbatim}

本题实现的 \texttt{swizzle\_\allowbreak{}128B} 会在 3.2 中直接用于
shared memory staging，并接受实际硬件上的 GEMM 判测。若 3.2 或 4.1
出现问题，可以先运行本题的判测，排除 swizzle 布局错误。

\begin{answer}

硬件：本地 AMD Ryzen AI 9 HX 370（CPU 实验，compute capability
不适用）；nvcc 12.9.86。

令行内逻辑 byte 地址为 \(c\)，每个 16B chunk 的低四位保持不动：

\[\begin{aligned}
p_{128}&=128r+(c\oplus((r\mathbin{\&}7)\ll4)),\\
p_{64}&=64r+(c\oplus((r\mathbin{\&}3)\ll4)),\\
p_{32}&=32r+(c\oplus((r\mathbin{\&}1)\ll4)).
\end{aligned}\]

对应异或位段：128B 为地址
\texttt{{[}6:\allowbreak{}4{]}\ \^{}=\allowbreak{}\ {[}9:\allowbreak{}7{]}}，64B
为
\texttt{{[}5:\allowbreak{}4{]}\ \^{}=\allowbreak{}\ {[}7:\allowbreak{}6{]}}，32B
为 \texttt{{[}4{]}\ \^{}=\allowbreak{}\ {[}5{]}}。固定行内 XOR
自反，行号不变，所以映射双射；固定逻辑 chunk 时，8/4/2
行分别遍历所有物理 chunk，因此满足题目的列访问判测。

\texttt{ARCH=\allowbreak{}89\ make\ -\allowbreak{}B\ run/\allowbreak{}m2\_\allowbreak{}smem/\allowbreak{}03\_\allowbreak{}swizzle}
输出 \texttt{128B\ PASS}、\texttt{64B\ PASS}、\texttt{32B\ PASS}。3.2
也已通过硬件 GEMM 对拍，进一步验证 128B staging 与 descriptor 配对。

\end{answer}

\begin{lookback}

课件 C12--C14 给出了使用 wgmma 实现单 tile
的完整示例，本作业不单独设置对应题目。wgmma 是
\texttt{sm\_\allowbreak{}90a} 专属指令，5090（sm120）和
B300（sm100）均无法运行；有兴趣的同学可以在集群 H 卡上自行尝试。

本模块重点保留 wgmma 路径中可以继续使用的两个部分：descriptor 与
swizzle。2.2、2.3 先分别完成它们的推导与判测，后面的 M3
会在实际硬件上继续使用。

\end{lookback}

\section{sm100：tcgen05}\label{sm100tcgen05}

在 sm100 上，tcgen05 将累加器放入 TMEM，并使用 mbarrier
处理异步完成通知。本模块会先在 B300 上完成一个单 tile
GEMM，再通过后续实验理解 TMEM、mbarrier 和 2-CTA MMA 的使用方式。

\begin{reading}

课件 S071--S093（TMEM：S073--S075；mma 与 idesc：S076--S079；
mbarrier：S080--S084；ld 与同步：S085--S086；七步流程：S087/F27；
2-CTA：S091）；PTX ISA 的 tcgen05 一族小节（alloc / mma / commit / ld /
fence）与 mbarrier 一节。

\end{reading}

\prob{3.1}{CONCEPT}

判断下列说法是否正确，并给出一句理由。其中 (d) 需要写出计算过程。

\begin{enumerate}
\def\labelenumi{(\alph{enumi})}
\item
  \texttt{tcgen05.\allowbreak{}ld} 读取 TMEM 时，每个 warp
  只能读取自己对应的 32 条 lane，warp 之间不能互相读取。
\item
  与 \texttt{mma.\allowbreak{}sync} 由 warp 协作、wgmma 由 warpgroup
  协作不同，\texttt{tcgen05.\allowbreak{}mma}
  由单个线程发射，随后由硬件异步执行。
\item
  TMEM 中的累加结果可以直接通过 TMA 搬回 global
  memory，不需要经过寄存器。
\item
  TMEM 每个 SM 包含 128 lane × 512 column × 4 B；一个 m128n256 的 f32
  accumulator 恰好占用其中一半。
\item
  \texttt{tcgen05.\allowbreak{}commit} 会阻塞直到之前发射的 mma
  全部完成，因此 commit 返回后即可安全读取 TMEM。
\end{enumerate}

\begin{answer}

硬件：NVIDIA B300 SXM6 AC，compute capability 10.3，148 SM；CUDA 13.0 /
nvcc 13.0.88，\texttt{ARCH=\allowbreak{}100f}。

\begin{enumerate}
\def\labelenumi{(\alph{enumi})}
\item
  对，本题使用的 TMEM load warp 布局中，每个 warp 访问对应的 32 行，因此
  128 行输出要由四个 warp 合作读回。
\item
  对，\texttt{tcgen05.\allowbreak{}mma} 由选出的单线程发射；alloc/ld
  等指令按各自的 warp 一致性要求执行。
\item
  错，结果通过
  \texttt{tcgen05.\allowbreak{}ld\ →\ registers\ →\ global\ store}
  写回。
\item
  对，总容量 \(128\times512\times4=262144\) B = 256 KiB；m128n256
  accumulator 是 \(128\times256\times4=131072\) B = 128 KiB，恰好一半。
\item
  错，commit 发起异步完成通知，必须等待对应 mbarrier phase 完成，再执行
  \texttt{tcgen05.\allowbreak{}fence:\allowbreak{}:\allowbreak{}after\_\allowbreak{}thread\_\allowbreak{}sync}
  和 TMEM load。
\end{enumerate}

\end{answer}

\begin{capstone}{prob 3.2(FROM-SCRATCH):tcgen05 单 tile GEMM \hfill {\footnotesize\ttfamily cuda/m3\_tcgen05/02\_single\_tile.cu}}

从零实现一个 tcgen05 单 tile GEMM：计算 m128n64k64 的 bf16 矩阵乘，使用
f32
累加、\texttt{cta\_\allowbreak{}group:\allowbreak{}:\allowbreak{}1}，并由一个
128 线程的 block 完成。

数据按照下面的路径流动：

\texttt{global\ →\ smem（K-\allowbreak{}major\ +\ 128B\ swizzle）→\ tcgen05.\allowbreak{}mma\ →\ TMEM\ →\ tcgen05.\allowbreak{}ld\ →\ global}

代码骨架只保留课件 F27 中的七步流程注释。descriptor 使用 2.2
中实现的编码方式，shared memory staging 使用 2.3 中的
\texttt{swizzle\_\allowbreak{}128B}。TMEM alloc、mbarrier、idesc 以及
mma / ld 的具体写法需要根据 PTX ISA 完成，相关语义可参考课件 C15--C21。

\texttt{fence.\allowbreak{}proxy.\allowbreak{}async} 的位置以及
\texttt{tcgen05.\allowbreak{}ld} 的 warp 可见范围已经在文件头中给出。

判测：

\begin{verbatim}
cd assignment02/cuda
make run/m3_tcgen05/02_single_tile
cd m3_tcgen05 && ./judge_tile.sh
\end{verbatim}

在正确版本通过后，故意去掉
\texttt{fence.\allowbreak{}proxy.\allowbreak{}async}
再运行一次，记录出现的现象并写入报告。结合 2.1(a) 中的排序问题解释原因。

\end{capstone}

\begin{answer}

硬件：NVIDIA B300 SXM6 AC，compute capability 10.3，148 SM；CUDA 13.0 /
nvcc 13.0.88，\texttt{ARCH=\allowbreak{}100f}。

完成 128-thread CTA 的七步路径：初始化 count=1 的 barrier、warp 0 分配
64 列 TMEM、全 CTA swizzled staging、proxy fence、单线程发四条 k16 MMA
并 commit、等待并由四个 warp 回读、释放 TMEM。首条 MMA 的
input-D=false，其后三条累加；load 后执行
\texttt{tcgen05.\allowbreak{}wait:\allowbreak{}:\allowbreak{}ld}。

\begin{Shaded}
\begin{Highlighting}[]
\BuiltInTok{cd}\NormalTok{ assignment02/cuda}
\FunctionTok{make} \AttributeTok{{-}B}\NormalTok{ run/m3\_tcgen05/02\_single\_tile}
\KeywordTok{(}\BuiltInTok{cd}\NormalTok{ m3\_tcgen05 }\KeywordTok{\&\&} \FunctionTok{bash}\NormalTok{ judge\_tile.sh}\KeywordTok{)}
\end{Highlighting}
\end{Shaded}

五个 seed 全部
\texttt{PASS}、\texttt{JUDGE:\allowbreak{}\ PASS}（\texttt{results/\allowbreak{}b300/\allowbreak{}judge-\allowbreak{}tcgen.\allowbreak{}log}）。

通过
\texttt{-\allowbreak{}DOMIT\_\allowbreak{}PROXY\_\allowbreak{}FENCE}
构建省略 proxy fence 的版本，运行相同多 seed 判测，五个 seed 全部
PASS，见
\texttt{results/\allowbreak{}b300/\allowbreak{}no-\allowbreak{}fence.\allowbreak{}log}。该版本
generic 写与 async 读的可见性依赖硬件时序；规范要求用 proxy fence
建立排序。最终版本保留 fence，并通过多 seed 判测。

\end{answer}

\prob{3.3}{DEBUG}[cuda/m3\_tcgen05/03\_bug\_mbarrier.cu]

这个程序连续发射多轮 mma，并在每一轮后读取结果，用来模拟 M4 中沿 K
维流式计算的过程。当前程序的 mbarrier
使用存在错误，判测带有超时机制，因此程序挂死本身也是需要观察的现象。

先运行：

\begin{verbatim}
cd assignment02/cuda
make bin/m3_tcgen05/03_bug_mbarrier
cd m3_tcgen05 && ./judge_mbar.sh
\end{verbatim}

\begin{enumerate}
\def\labelenumi{(\alph{enumi})}
\item
  分别记录 \texttt{rounds\ =\allowbreak{}\ 1}、\texttt{2}、\texttt{4}
  时的运行结果以及问题出现的条件。
\item
  修复程序，并画出修改前后 mbarrier 的状态变化，包括 phase 和 arrival
  count。指出错误版本在哪一次等待中过早或过晚放行，并解释为什么会产生
  (a) 中观察到的现象。
\end{enumerate}

提示：注意 \texttt{tcgen05.\allowbreak{}ld} 与尚未完成的 mma
之间的关系。

\begin{answer}

硬件：NVIDIA B300 SXM6 AC，compute capability 10.3，148 SM；CUDA 13.0 /
nvcc 13.0.88，\texttt{ARCH=\allowbreak{}100f}。

原始版本 seed=42：rounds=1 为 \texttt{PASS}；rounds=2、4 都在 15
秒超时后退出 124。修复后 rounds=1、2、4 × seeds=42、7 全部
\texttt{PASS}，最终 \texttt{JUDGE:\allowbreak{}\ PASS}。

正确状态机（count 每轮重新装载为 1）：

\begin{Shaded}
\begin{Highlighting}[]
\NormalTok{初始 (phase=0, pending=1)}
\NormalTok{round 0: commit 完成 → (0,0) → (1,1); wait(0) 放行}
\NormalTok{round 1: commit 完成 → (1,0) → (0,1); wait(1) 放行}
\NormalTok{round 2: commit 完成 → (0,0) → (1,1); wait(0) 放行}
\NormalTok{round 3: commit 完成 → (1,0) → (0,1); wait(1) 放行}
\end{Highlighting}
\end{Shaded}

错误版本每次都 \texttt{wait(0)}。第一轮后 phase=1，第二轮 wait(0)
可能在第二条 MMA 完成前过早放行，使 TMEM load 与在飞 MMA 冲突；若第二轮
MMA 已完成，phase 回到 0，wait(0)
会等待下一次翻转，此时下一轮发射也被阻塞，形成死锁。具体表现取决于 MMA
完成与 wait 检查的先后顺序。

修复为
\texttt{mbar\_\allowbreak{}wait(mbar\_\allowbreak{}u32,\allowbreak{}\ round\ \&\ 1)}；读完后用
\texttt{wait:\allowbreak{}:\allowbreak{}ld}、\texttt{fence:\allowbreak{}:\allowbreak{}before\_\allowbreak{}thread\_\allowbreak{}sync}
和 CTA barrier，保证下一轮写 TMEM 前四个 warp
均完成读取。源码快照与原始日志分别在
\texttt{results/\allowbreak{}original/\allowbreak{}03\_\allowbreak{}bug\_\allowbreak{}mbarrier.\allowbreak{}cu}、\texttt{results/\allowbreak{}b300/\allowbreak{}debug-\allowbreak{}before.\allowbreak{}log}。

\end{answer}

\prob{3.4}{EXPERIMENT}[cuda/m3\_tcgen05/04\_cta\_pair.cu]

比较 \texttt{cta\_\allowbreak{}group:\allowbreak{}:\allowbreak{}1} 和
\texttt{cta\_\allowbreak{}group:\allowbreak{}:\allowbreak{}2} 完成同一个
m256n64k64 任务时的数据开销。

\texttt{cta\_\allowbreak{}group:\allowbreak{}:\allowbreak{}1}
使用两个独立
block；\texttt{cta\_\allowbreak{}group:\allowbreak{}:\allowbreak{}2}
使用一个 cluster 发射一条 M=256 的 mma，两个 CTA 分别保存 B
矩阵的一部分。

先预测，再运行实验：

\begin{enumerate}
\def\labelenumi{(\alph{enumi})}
\item
  在 \texttt{cta\_\allowbreak{}group:\allowbreak{}:\allowbreak{}2}
  中，每个 CTA 所需的 B shared memory 是
  \texttt{cta\_\allowbreak{}group:\allowbreak{}:\allowbreak{}1}
  的多少？TMEM 的占用又如何变化？程序会打印 shared memory
  用量，用它验证你的预测。
\item
  使用 Nsight Compute 比较两种实现的 shared memory 总流量。
\item
  \texttt{cta\_\allowbreak{}group:\allowbreak{}:\allowbreak{}2}
  节省下来的 shared memory 容量，在 M4 的 pipeline 中可以用来做什么？
\item
  \texttt{cta\_\allowbreak{}group:\allowbreak{}:\allowbreak{}2} 依赖
  sm90 引入的哪一种硬件机制？5090 不支持 2-CTA MMA，结合 0.2
  中整理的硬件参数，说明为什么这类机制更常出现在数据中心 GPU 上。
\end{enumerate}

运行：

\begin{verbatim}
cd assignment02/cuda
make run/m3_tcgen05/04_cta_pair
\end{verbatim}

在这个单 tile
实验中，两种实现的运行时间差异处于噪声范围内，因此不要用耗时判断优劣。主要比较
shared memory 用量以及 Nsight Compute 测得的流量。

\begin{answer}

硬件：NVIDIA B300 SXM6 AC，compute capability 10.3，148 SM；CUDA 13.0 /
nvcc 13.0.88，\texttt{ARCH=\allowbreak{}100f}。

\texttt{make\ -\allowbreak{}B\ run/\allowbreak{}m3\_\allowbreak{}tcgen05/\allowbreak{}04\_\allowbreak{}cta\_\allowbreak{}pair}
两条路径均 \texttt{PASS(bad=\allowbreak{}0)}。

\begin{enumerate}
\def\labelenumi{(\alph{enumi})}
\item
  group 1 每 CTA 保存 B 的 \(64\times64\times2=8192\) B，group 2 保存
  \(32\times64\times2=4096\) B，即一半。A 仍为 16384 B。程序打印总
  shared/CTA 为 24588 B 和 20492 B（含 barrier/地址等静态对象），减少
  4096 B。每 CTA 的输出为 128×64 f32，TMEM 申请 64 列即 32 KiB；两 CTA
  总计 64 KiB。
\item
  用 NCU 记录前两个 launch，LSU shared store wavefront 总量由 778 降为
  650，减少 128，约 16.45\%；load 为 26 和 29，额外 cluster
  控制访问会影响这一计数。两个 CTA 的显式矩阵 staging 字节量从 49152 B
  减为 40960 B，减少 16.67\%，与 store 指标一致。这里的 LSU 指标统计普通
  load/store 路径的 shared 访问。
\item
  每个 stage 节省的 B 空间可用于增加缓冲级数或扩大输出
  tile，代价是需配合 cluster 同步和两个 SM 的资源。
\item
  依赖 Thread Block Cluster、跨 CTA shared memory
  通信及协同驻留。数据中心 GPU 的宽 Tensor Core、HBM/NVLink
  和大型训练任务更容易摊薄这种协同成本；2-CTA tcgen05 属于 sm100 系列的
  MMA 能力。
\end{enumerate}

实验按 shared 用量和 LSU 计数比较两种实现。原始用量/判测见
\texttt{04\_\allowbreak{}cta\_\allowbreak{}pair.\allowbreak{}log}，NCU
命令与数据见 \texttt{ncu-\allowbreak{}cta.\allowbreak{}csv}。

\end{answer}

\section{完整 GEMM}\label{ux5b8cux6574-gemm}

本模块将 3.2 的单 tile 实现扩展为完整的 B300 GEMM，并依次加入
tiling、TMA 和多级 pipeline。实验统一使用 4096³ 的 bf16 GEMM， tile
大小固定为 128×64×64，并与 cuBLAS 结果进行严格相等比较。

\begin{reading}

课件 S084（pipeline 骨架）、S087；Session 4 讲义对应章节；PTX ISA 的
\texttt{cp.\allowbreak{}async.\allowbreak{}bulk.\allowbreak{}tensor}、mbarrier（\texttt{expect\_\allowbreak{}tx}）小节；CUDA
Driver API 的 \texttt{cuTensorMapEncodeTiled}。

\end{reading}

下面的表格贯穿 4.1--4.3。第 0 行需要在同一块 GPU 上重新运行 assignment01
Bonus 中的 naive matmul。由于该实现使用 fp32，只比较性能量级。

\begin{longtable}[]{@{}llll@{}}
\toprule\noalign{}
实现 & TFLOPS & 对 cuBLAS 达成率 & 一句话：时间主要花在哪 \\
\midrule\noalign{}
\endhead
\bottomrule\noalign{}
\endlastfoot
naive（assignment01，fp32） & & & \\
4.1 tiled & & & \\
4.2 TMA & & & \\
4.3 pipeline（S=3） & & & \\
cuBLAS & & 100\% & \\
\end{longtable}

\prob{4.1}{FROM-SCRATCH}[cuda/m4\_gemm/01\_tiled.cu]

将 3.2 的单 tile 实现扩展为完整 GEMM：使用 grid 覆盖所有输出 tile， 并沿
K 维循环完成累加。数据 staging 仍然使用 \texttt{st.\allowbreak{}shared}
和 swizzle。

相对 3.2，本题主要新增 tile 映射以及 K 循环中的同步和累加。
具体步骤见文件头。

\begin{verbatim}
cd assignment02/cuda
make run/m4_gemm/01_tiled
\end{verbatim}

通过判测后，填写性能表中 \texttt{4.\allowbreak{}1\ tiled} 一行。结合 0.2
中计算的机器 平衡点，判断此时性能主要受哪一环节限制。

\begin{answer}

硬件：NVIDIA B300 SXM6 AC，compute capability 10.3，148 SM；CUDA 13.0 /
nvcc 13.0.88，\texttt{ARCH=\allowbreak{}100f}。

\texttt{make\ -\allowbreak{}B\ run/\allowbreak{}m4\_\allowbreak{}gemm/\allowbreak{}01\_\allowbreak{}tiled}：4096³
\texttt{PASS(bad=\allowbreak{}0)}，41.4 TFLOPS；cuBLAS BF16 为 1767.0
TFLOPS，相对达成率 2.34\%。见
\texttt{results/\allowbreak{}b300/\allowbreak{}01\_\allowbreak{}tiled.\allowbreak{}log}。

grid 为 32×64，K 分 64 轮。每轮全 CTA 进行 global load、swizzle
地址计算和 shared store，proxy fence 后由线程 0 发四条
MMA；等待本轮消费完成后才覆写 shared，accumulator 一直留在 TMEM。

每个 K tile 的理想 staging 强度为
\(2\cdot128\cdot64\cdot64/((128+64)\cdot64\cdot2)=42.67\)
FLOP/B，低于机器平衡点 281.25。NCU 实测 DRAM 0.31\%、SM throughput
19.07\%、active warp 43.80\%，主要瓶颈是指令化
staging、地址计算和同步供数。

\end{answer}

\prob{4.2}{MODIFY}[cuda/m4\_gemm/02\_tma.cu]

将 4.1 中的 shared memory staging 改为 TMA。

host 端使用 \texttt{cuTensorMapEncodeTiled} 创建 tensor map；kernel
中使用
\texttt{cp.\allowbreak{}async.\allowbreak{}bulk.\allowbreak{}tensor}
搭配 mbarrier 的 \texttt{expect\_\allowbreak{}tx} 完成搬运。
本题仍然使用单缓冲。

tensor map 的参数以及同步方式的变化已经在文件头中给出。 swizzle 由 TMA
根据 tensor map 自动完成，原有 descriptor 不需要修改字段。

\begin{verbatim}
cd assignment02/cuda
make run/m4_gemm/02_tma
\end{verbatim}

通过判测后，填写性能表中 \texttt{4.\allowbreak{}2\ TMA} 一行。

比较 4.1 与 4.2，并回答：4.1 中 shared memory staging 的开销由哪些
部分组成？改用 TMA 后，其中哪些工作不再由普通 CUDA 指令完成？

使用 Nsight Compute 辅助分析，观察 4.1 中 SM 时间主要消耗在哪些部分。

\begin{answer}

硬件：NVIDIA B300 SXM6 AC，compute capability 10.3，148 SM；CUDA 13.0 /
nvcc 13.0.88，\texttt{ARCH=\allowbreak{}100f}。

\texttt{make\ -\allowbreak{}B\ run/\allowbreak{}m4\_\allowbreak{}gemm/\allowbreak{}02\_\allowbreak{}tma}：4096³
\texttt{PASS(bad=\allowbreak{}0)}，568.5 TFLOPS，对 cuBLAS 1765.7
的达成率约 32.20\%；见 \texttt{02\_\allowbreak{}tma.\allowbreak{}log}。

A/B 的 tensor map 分别取
dims=\{K,M\}/\{K,N\}、stride=\{2K\}、box=\{64,128\}/\{64,64\}；检查
\texttt{cuTensorMapEncodeTiled} 返回值。每轮一次
\texttt{arrive.\allowbreak{}expect\_\allowbreak{}tx} 报 24576 B，两条
TMA copy 共用 full barrier；消费端等待 full 后发 MMA，commit 到
empty，下一轮搬运前等待 empty。

TMA 接管普通 load/store、逐元素地址计算和 swizzle，将 global
数据直接搬入 shared，原来的 SM100 MMA descriptor 保持不变。TMA 写和 MMA
读同在 async proxy，同步通过数据到达通知与 tcgen05 的线程排序完成。相对
4.1 提升约 13.73 倍，主要收益来自 staging 指令开销的减少。

\end{answer}

\begin{capstone}{prob 4.3(FROM-SCRATCH):多级流水 \hfill {\footnotesize\ttfamily cuda/m4\_gemm/03\_pipeline.cu}}

将 4.2 的单缓冲改为 \texttt{STAGES} 级循环缓冲，使后续 K tile 的 TMA
预取能够与当前 tile 的 mma 计算重叠。

\texttt{STAGES} 为编译参数，例如：

\begin{verbatim}
STAGES=4 make -B run/m4_gemm/03_pipeline
\end{verbatim}

这里的 \texttt{-\allowbreak{}B} 不能省略，否则修改 \texttt{STAGES}
后可能不会重新编译。

文件头给出了每个 stage 使用双 mbarrier 的基本结构，同时包含一个需要
处理的 pipeline hazard：强制发射与机会式预取之间的边界处理错误会导致
死锁。在该错误下，1024³ 可能偶尔通过，而 4096³ 会稳定暴露问题。
开始实现前先阅读文件头说明。

本题不要求 warp specialization、persistent kernel 或 epilogue 融合，
也不设置 cuBLAS 达成率门槛。重点是流水实现是否正确，以及能否根据实验
结果解释性能变化。

\begin{verbatim}
cd assignment02/cuda
make run/m4_gemm/03_pipeline
cd m4_gemm && ./sweep_stages.sh
\end{verbatim}

完成以下内容：

\begin{enumerate}
\def\labelenumi{\arabic{enumi}.}
\item
  使用 \texttt{S=\allowbreak{}3} 的结果填写性能表中
  \texttt{4.\allowbreak{}3\ pipeline} 一行。
\item
  完成不同 stage 数的性能测试：

  \begin{longtable}[]{@{}lllll@{}}
  \toprule\noalign{}
  形状 & S=2 & S=3 & S=4 & S=6 \\
  \midrule\noalign{}
  \endhead
  \bottomrule\noalign{}
  \endlastfoot
  4096³ & & & & \\
  256 × 4096 × 16384 & & & & \\
  \end{longtable}

  比较两个形状对 \texttt{STAGES} 的敏感程度，并结合 shared memory 用量、
  每个 SM 可同时驻留的 block 数以及 block 间并发能够隐藏的延迟进行解释。
\item
  任选一个 \texttt{STAGES}，画出稳态阶段各 stage 中 TMA 与 mma
  的流水时空图。
\item
  回答：

  \begin{enumerate}
  \def\labelenumii{(\alph{enumii})}
  \item
    从 4.1 到 4.3，主要瓶颈发生了哪些变化？
  \item
    梯子表中每一级优化分别减少了哪部分开销？
  \item
    如果继续增大 tile 或增加 stage 数，shared memory 与 TMEM
    哪一个会先成为容量限制？结合 3.4(c) 的结果说明。
  \end{enumerate}
\end{enumerate}

\end{capstone}

\begin{answer}

硬件：NVIDIA B300 SXM6 AC，compute capability 10.3，148 SM；CUDA 13.0 /
nvcc 13.0.88，\texttt{ARCH=\allowbreak{}100f}。

每 stage 使用独立 full/empty barrier 与循环缓冲。预热发 min(S,iters) 个
TMA；顺序记录 next，当前 it 待发时阻塞等 empty 并补发，深预取使用非阻塞
try-wait。full 的 parity 为 \texttt{(it/\allowbreak{}S)\&1}，复用时等待
empty 的上一代 parity 为
\texttt{((next/\allowbreak{}S)-\allowbreak{}1)\&1}；每轮先保证当前 copy
已发射，再等待 full，最终 drain 后读取 TMEM。

\begin{Shaded}
\begin{Highlighting}[]
\FunctionTok{make} \AttributeTok{{-}B}\NormalTok{ run/m4\_gemm/03\_pipeline}
\FunctionTok{bash}\NormalTok{ m4\_gemm/sweep\_stages.sh}
\end{Highlighting}
\end{Shaded}

所有 8 个形状/级数组合均 \texttt{PASS(bad=\allowbreak{}0)}；另测
S=3、M=128、N=64、K=64/128/192/256（预热未满、刚好填满、首次复用），四项均
PASS，见
\texttt{pipeline-\allowbreak{}short-\allowbreak{}k.\allowbreak{}log}。以下
TFLOPS 取
\texttt{results/\allowbreak{}b300/\allowbreak{}stages.\allowbreak{}log}：

\begin{longtable}[]{@{}lrrrr@{}}
\toprule\noalign{}
形状 & S=2 & S=3 & S=4 & S=6 \\
\midrule\noalign{}
\endhead
\bottomrule\noalign{}
\endlastfoot
4096³ & 573.3 & 574.5 & 468.5 & 281.9 \\
256×4096×16384 & 212.2 & 271.1 & 271.9 & 275.1 \\
\end{longtable}

shared 主体每 stage 为 24 KiB，动态分配另加 1024B 对齐余量；S=2/3/4/6
分别是 50176/74752/99328/148480 B，静态元数据为 48/64/80/112
B，寄存器均为 32/thread。按本机 233472 B shared/SM 计算，shared
容量对应的驻留上界依次为 4/3/2/1 CTA；实际驻留由
shared、寄存器及其他硬件资源共同决定。占用 API 对本 tcgen05 kernel 返回
1。NCU 的 active-warp 实测依次为 21.03\%、16.82\%、12.00\%、6.23\%，呈现
shared 用量增加、活跃 warp
比例下降的趋势（\texttt{ncu-\allowbreak{}stage-\allowbreak{}*.\allowbreak{}csv}）。

大 grid 有 2048 CTA，stage 加深带来的 shared 占用会削弱块间并发；小 grid
仅 128 CTA，小于 148 SM，已有并发不足，长 K 的等待更依赖块内预取，因此
S=2→3 的收益显著，继续加深逐渐饱和。

S=3 的稳态示意（箭头表示先等该 stage 的
empty，再复用；非等比例时间轴）：

\begin{Shaded}
\begin{Highlighting}[]
\NormalTok{时间 →      t0           t1           t2           t3}
\NormalTok{stage 0   MMA(k0) ──→ TMA(k3) ────────────────→ MMA(k3)}
\NormalTok{stage 1   TMA(k1) ──→ MMA(k1) ──→ TMA(k4)}
\NormalTok{stage 2   TMA(k2) ────────────→ MMA(k2) ──→ TMA(k5)}
\end{Highlighting}
\end{Shaded}

梯子表（4096³，同一 B300；naive 是 assignment01 的原
kernel，仅换同形状计时和 cuBLAS reference）：

\begin{longtable}[]{@{}
  >{\raggedright\arraybackslash}p{(\linewidth - 6\tabcolsep) * \real{0.2143}}
  >{\raggedleft\arraybackslash}p{(\linewidth - 6\tabcolsep) * \real{0.2857}}
  >{\raggedleft\arraybackslash}p{(\linewidth - 6\tabcolsep) * \real{0.2857}}
  >{\raggedright\arraybackslash}p{(\linewidth - 6\tabcolsep) * \real{0.2143}}@{}}
\toprule\noalign{}
\begin{minipage}[b]{\linewidth}\raggedright
实现
\end{minipage} & \begin{minipage}[b]{\linewidth}\raggedleft
TFLOPS
\end{minipage} & \begin{minipage}[b]{\linewidth}\raggedleft
对 cuBLAS BF16 达成率
\end{minipage} & \begin{minipage}[b]{\linewidth}\raggedright
主要剩余开销
\end{minipage} \\
\midrule\noalign{}
\endhead
\bottomrule\noalign{}
\endlastfoot
naive FP32 & 6.4 & 0.36\% & 标量 FMA 与重复 global 访问 \\
tiled & 41.4 & 2.34\% & 指令化 staging、swizzle、等待 \\
TMA & 568.5 & 32.17\% & 单缓冲搬运与计算串行 \\
pipeline S=3 & 576.3 & 32.61\% & 固定窄 tile
的供数/发射/同步与驻留限制 \\
cuBLAS BF16 & 1767.0 & 100\% & 库内部优化基线 \\
\end{longtable}

naive 的 FP32 cuBLAS 对照为 65.7 TFLOPS。梯子表统一以 BF16 cuBLAS
为分母比较性能量级，naive 使用 FP32，其余实现使用 BF16 输入。4.3 相对
4.2 小幅提高，剩余开销集中于固定窄 tile 的供数、发射、同步与驻留。

固定输出 tile 时，S 增大线性增加 shared，TMEM 始终为 128×64×4=32768
B，因此先达到 shared 容量限制。加大 M/N 会同时增加 TMEM 用量；2-CTA
通过节省 B shared 为更多 stage 留出空间，同时需要考虑 cluster 同步和
TMEM 容量。

\end{answer}

\prob[Optional]{4.4}{FROM-SCRATCH}

任选一个方向继续优化：

\begin{enumerate}
\def\labelenumi{(\alph{enumi})}
\item
  实现 4.3 的
  \texttt{cta\_\allowbreak{}group:\allowbreak{}:\allowbreak{}2}
  版本，利用 B shared memory 用量的减少 增加 pipeline 深度或扩大 tile；
\item
  在 5090 上使用 \texttt{mma.\allowbreak{}sync} 实现相同的 tiling，并与
  B300 的结果比较；
\item
  自由优化当前 kernel，提高相对 cuBLAS 的性能，并记录每一步优化
  解决了什么问题。
\end{enumerate}

\begin{answer}

硬件：NVIDIA B300 SXM6 AC，compute capability 10.3，148 SM；CUDA 13.0 /
nvcc 13.0.88，\texttt{ARCH=\allowbreak{}100f}。

本次完成必做主线，未选择 4.4 的额外 GEMM 优化方向。5.4
另外完成了针对融合 kernel 访存的优化与复测。

\end{answer}

\prob{4.5}{EXPERIMENT}[cuda/m4\_gemm/05\_thin\_gemm.cu]

观察矩阵形状较窄时 Tensor Core GEMM 的性能变化。

实验形状取自 vLLM 主树中 Kimi K3 的 decode GEMM dispatch 表：

\texttt{vllm/\allowbreak{}models/\allowbreak{}kimi\_\allowbreak{}k3/\allowbreak{}nvidia/\allowbreak{}low\_\allowbreak{}latency\_\allowbreak{}gemm.\allowbreak{}py}

使用其中七个投影层对应的 N、K，并对 M 取：

\(\{1, 8, 16, 64, 256, 1024, 4096, 16384, 65536\}\)

其中 N 和 K 由权重形状决定，变化的 M 表示一次 GEMM 处理的 token 数量。
较小的 M 对应 decode batch；较大的 M 可以对应 chunked prefill 中单次
处理的 chunk。vLLM 在 \(M \le 16\) 时会放弃 cuBLAS / Tensor Core 路径，
改用 CUDA Core FMA 的 skinny kernel。

运行程序前，先对每个形状计算 arithmetic intensity：

\[
AI = \frac{2MNK}{2MK + 2NK + 2MN}
\]

将结果与 0.2 中得到的机器平衡点比较，并计算对应的理论性能上限：

\begin{itemize}
\tightlist
\item
  compute bound：Tensor Core 峰值；
\item
  memory bound：\(AI \times\) 显存带宽。
\end{itemize}

然后运行：

\begin{verbatim}
cd assignment02/cuda
make bin/m4_gemm/05_thin_gemm
./bin/m4_gemm/05_thin_gemm <峰值TFLOPS> <带宽GB/s>
\end{verbatim}

其中两个参数使用 0.2 中得到的数值。

程序会输出 TFLOPS、GB/s、AI 以及相对于 compute roof 和 memory roof
的达成率。根据结果回答：

\begin{enumerate}
\def\labelenumi{(\alph{enumi})}
\item
  描述性能随 M 变化的趋势。M 较小时，Tensor Core 性能从什么时候
  开始明显下降？M
  增大到什么范围后进入相对稳定的平台？平台上的达成率是多少？
\item
  对性能下降的形状，比较 compute roof 和 memory roof 两种达成率。
  哪些形状主要受到显存带宽限制？
\item
  \texttt{f\_\allowbreak{}b\_\allowbreak{}proj}（K=128）中两个 roof
  的达成率都较低。结合矩阵形状分析， 它的限制来自哪里？
\item
  根据上述结果解释 vLLM 在 \(M \le 16\) 时选择 skinny CUDA Core kernel
  的原因。
\end{enumerate}

\texttt{in\_\allowbreak{}proj\_\allowbreak{}qkvgfab} 对应 KDA
的输入投影。完成团队题 C1/C2 的同学可以
使用这一行的实验结果作为后续分析的参考。

\begin{answer}

硬件：NVIDIA B300 SXM6 AC，compute capability 10.3，148 SM；CUDA 13.0 /
nvcc 13.0.88，\texttt{ARCH=\allowbreak{}100f}。

运行
\texttt{.\allowbreak{}/\allowbreak{}bin/\allowbreak{}m4\_\allowbreak{}gemm/\allowbreak{}05\_\allowbreak{}thin\_\allowbreak{}gemm\ 2250\ 8000}；七个投影、九个
M 的完整原始输出见
\texttt{results/\allowbreak{}b300/\allowbreak{}thin.\allowbreak{}log}。AI
用题设 BF16 输入/输出字节口径，roof 为 \(\min(2250,8AI)\)
TFLOPS；compute 和 memory 达成率分别为实测/2250 与实测/(8AI)。

\end{answer}

\begin{longtable}[]{@{}
  >{\raggedright\arraybackslash}p{(\linewidth - 12\tabcolsep) * \real{0.1111}}
  >{\raggedleft\arraybackslash}p{(\linewidth - 12\tabcolsep) * \real{0.1481}}
  >{\raggedleft\arraybackslash}p{(\linewidth - 12\tabcolsep) * \real{0.1481}}
  >{\raggedleft\arraybackslash}p{(\linewidth - 12\tabcolsep) * \real{0.1481}}
  >{\raggedleft\arraybackslash}p{(\linewidth - 12\tabcolsep) * \real{0.1481}}
  >{\raggedleft\arraybackslash}p{(\linewidth - 12\tabcolsep) * \real{0.1481}}
  >{\raggedleft\arraybackslash}p{(\linewidth - 12\tabcolsep) * \real{0.1481}}@{}}
\toprule\noalign{}
\begin{minipage}[b]{\linewidth}\raggedright
投影
\end{minipage} & \begin{minipage}[b]{\linewidth}\raggedleft
M
\end{minipage} & \begin{minipage}[b]{\linewidth}\raggedleft
AI
\end{minipage} & \begin{minipage}[b]{\linewidth}\raggedleft
理论 roof TFLOPS
\end{minipage} & \begin{minipage}[b]{\linewidth}\raggedleft
实测 TFLOPS
\end{minipage} & \begin{minipage}[b]{\linewidth}\raggedleft
compute roof 达成率
\end{minipage} & \begin{minipage}[b]{\linewidth}\raggedleft
memory roof 达成率
\end{minipage} \\
\midrule\noalign{}
\endhead
\bottomrule\noalign{}
\endlastfoot
\texttt{f\_\allowbreak{}b\_\allowbreak{}proj} & 1 & 0.99 & 7.93 & 0.1 &
0.0\% & 1.3\% \\
\texttt{f\_\allowbreak{}b\_\allowbreak{}proj} & 8 & 7.49 & 59.94 & 0.8 &
0.0\% & 1.3\% \\
\texttt{f\_\allowbreak{}b\_\allowbreak{}proj} & 16 & 14.09 & 112.73 &
1.6 & 0.1\% & 1.4\% \\
\texttt{f\_\allowbreak{}b\_\allowbreak{}proj} & 64 & 41.51 & 332.11 &
6.6 & 0.3\% & 2.0\% \\
\texttt{f\_\allowbreak{}b\_\allowbreak{}proj} & 256 & 80.84 & 646.74 &
25.9 & 1.1\% & 4.0\% \\
\texttt{f\_\allowbreak{}b\_\allowbreak{}proj} & 1024 & 105.93 & 847.45 &
100.8 & 4.5\% & 11.9\% \\
\texttt{f\_\allowbreak{}b\_\allowbreak{}proj} & 4096 & 114.84 & 918.73 &
325.0 & 14.4\% & 35.4\% \\
\texttt{f\_\allowbreak{}b\_\allowbreak{}proj} & 16384 & 117.31 & 938.46
& 572.2 & 25.4\% & 61.0\% \\
\texttt{f\_\allowbreak{}b\_\allowbreak{}proj} & 65536 & 117.94 & 943.53
& 694.1 & 30.8\% & 73.6\% \\
\texttt{q\_\allowbreak{}b\_\allowbreak{}proj} & 1 & 1.00 & 7.99 & 0.8 &
0.0\% & 9.5\% \\
\texttt{q\_\allowbreak{}b\_\allowbreak{}proj} & 8 & 7.93 & 63.45 & 9.6 &
0.4\% & 15.1\% \\
\texttt{q\_\allowbreak{}b\_\allowbreak{}proj} & 16 & 15.73 & 125.82 &
31.5 & 1.4\% & 25.0\% \\
\texttt{q\_\allowbreak{}b\_\allowbreak{}proj} & 64 & 59.84 & 478.75 &
123.9 & 5.5\% & 25.9\% \\
\texttt{q\_\allowbreak{}b\_\allowbreak{}proj} & 256 & 200.35 & 1602.78 &
470.6 & 20.9\% & 29.4\% \\
\texttt{q\_\allowbreak{}b\_\allowbreak{}proj} & 1024 & 485.05 & 2250.00
& 1212.9 & 53.9\% & 31.3\% \\
\texttt{q\_\allowbreak{}b\_\allowbreak{}proj} & 4096 & 752.33 & 2250.00
& 1743.1 & 77.5\% & 29.0\% \\
\texttt{q\_\allowbreak{}b\_\allowbreak{}proj} & 16384 & 872.52 & 2250.00
& 1904.6 & 84.6\% & 27.3\% \\
\texttt{q\_\allowbreak{}b\_\allowbreak{}proj} & 65536 & 908.82 & 2250.00
& 2026.7 & 90.1\% & 27.9\% \\
\texttt{o\_\allowbreak{}proj} & 1 & 1.00 & 7.99 & 4.6 & 0.2\% &
58.2\% \\
\texttt{o\_\allowbreak{}proj} & 8 & 7.95 & 63.60 & 48.4 & 2.2\% &
76.1\% \\
\texttt{o\_\allowbreak{}proj} & 16 & 15.80 & 126.40 & 94.5 & 4.2\% &
74.8\% \\
\texttt{o\_\allowbreak{}proj} & 64 & 60.92 & 487.34 & 313.5 & 13.9\% &
64.3\% \\
\texttt{o\_\allowbreak{}proj} & 256 & 212.91 & 1703.29 & 930.1 & 41.3\%
& 54.6\% \\
\texttt{o\_\allowbreak{}proj} & 1024 & 565.89 & 2250.00 & 1366.0 &
60.7\% & 30.2\% \\
\texttt{o\_\allowbreak{}proj} & 4096 & 966.47 & 2250.00 & 1906.5 &
84.7\% & 24.7\% \\
\texttt{o\_\allowbreak{}proj} & 16384 & 1174.28 & 2250.00 & 2005.2 &
89.1\% & 21.3\% \\
\texttt{o\_\allowbreak{}proj} & 65536 & 1240.99 & 2250.00 & 2090.7 &
92.9\% & 21.1\% \\
\texttt{fused\_\allowbreak{}qkv\_\allowbreak{}a\_\allowbreak{}proj} & 1
& 1.00 & 8.00 & 3.9 & 0.2\% & 48.5\% \\
\texttt{fused\_\allowbreak{}qkv\_\allowbreak{}a\_\allowbreak{}proj} & 8
& 7.96 & 63.69 & 36.8 & 1.6\% & 57.9\% \\
\texttt{fused\_\allowbreak{}qkv\_\allowbreak{}a\_\allowbreak{}proj} & 16
& 15.84 & 126.76 & 72.4 & 3.2\% & 57.2\% \\
\texttt{fused\_\allowbreak{}qkv\_\allowbreak{}a\_\allowbreak{}proj} & 64
& 61.58 & 492.67 & 212.7 & 9.5\% & 43.2\% \\
\texttt{fused\_\allowbreak{}qkv\_\allowbreak{}a\_\allowbreak{}proj} &
256 & 221.28 & 1770.21 & 729.5 & 32.4\% & 41.2\% \\
\texttt{fused\_\allowbreak{}qkv\_\allowbreak{}a\_\allowbreak{}proj} &
1024 & 629.11 & 2250.00 & 1504.5 & 66.9\% & 29.9\% \\
\texttt{fused\_\allowbreak{}qkv\_\allowbreak{}a\_\allowbreak{}proj} &
4096 & 1166.68 & 2250.00 & 1754.7 & 78.0\% & 18.8\% \\
\texttt{fused\_\allowbreak{}qkv\_\allowbreak{}a\_\allowbreak{}proj} &
16384 & 1483.62 & 2250.00 & 1826.3 & 81.2\% & 15.4\% \\
\texttt{fused\_\allowbreak{}qkv\_\allowbreak{}a\_\allowbreak{}proj} &
65536 & 1591.72 & 2250.00 & 1737.2 & 77.2\% & 13.6\% \\
\texttt{in\_\allowbreak{}proj\_\allowbreak{}qkvgfab} & 1 & 1.00 & 8.00 &
5.3 & 0.2\% & 66.8\% \\
\texttt{in\_\allowbreak{}proj\_\allowbreak{}qkvgfab} & 8 & 7.98 & 63.85
& 45.8 & 2.0\% & 71.8\% \\
\texttt{in\_\allowbreak{}proj\_\allowbreak{}qkvgfab} & 16 & 15.92 &
127.39 & 86.6 & 3.9\% & 68.0\% \\
\texttt{in\_\allowbreak{}proj\_\allowbreak{}qkvgfab} & 64 & 62.80 &
502.40 & 383.1 & 17.0\% & 76.3\% \\
\texttt{in\_\allowbreak{}proj\_\allowbreak{}qkvgfab} & 256 & 237.82 &
1902.59 & 1007.8 & 44.8\% & 53.0\% \\
\texttt{in\_\allowbreak{}proj\_\allowbreak{}qkvgfab} & 1024 & 784.25 &
2250.00 & 1408.9 & 62.6\% & 22.5\% \\
\texttt{in\_\allowbreak{}proj\_\allowbreak{}qkvgfab} & 4096 & 1842.70 &
2250.00 & 1768.0 & 78.6\% & 12.0\% \\
\texttt{in\_\allowbreak{}proj\_\allowbreak{}qkvgfab} & 16384 & 2781.04 &
2250.00 & 1759.8 & 78.2\% & 7.9\% \\
\texttt{in\_\allowbreak{}proj\_\allowbreak{}qkvgfab} & 65536 & 3186.74 &
2250.00 & 1684.2 & 74.9\% & 6.6\% \\
\texttt{dense\_\allowbreak{}down\_\allowbreak{}proj} & 1 & 1.00 & 8.00 &
4.8 & 0.2\% & 60.4\% \\
\texttt{dense\_\allowbreak{}down\_\allowbreak{}proj} & 8 & 7.98 & 63.87
& 41.3 & 1.8\% & 64.7\% \\
\texttt{dense\_\allowbreak{}down\_\allowbreak{}proj} & 16 & 15.93 &
127.47 & 81.2 & 3.6\% & 63.7\% \\
\texttt{dense\_\allowbreak{}down\_\allowbreak{}proj} & 64 & 62.96 &
503.69 & 280.2 & 12.5\% & 55.6\% \\
\texttt{dense\_\allowbreak{}down\_\allowbreak{}proj} & 256 & 240.15 &
1921.17 & 1027.0 & 45.6\% & 53.5\% \\
\texttt{dense\_\allowbreak{}down\_\allowbreak{}proj} & 1024 & 810.08 &
2250.00 & 1208.9 & 53.7\% & 18.7\% \\
\texttt{dense\_\allowbreak{}down\_\allowbreak{}proj} & 4096 & 1991.95 &
2250.00 & 1810.4 & 80.5\% & 11.4\% \\
\texttt{dense\_\allowbreak{}down\_\allowbreak{}proj} & 16384 & 3135.63 &
2250.00 & 1667.5 & 74.1\% & 6.6\% \\
\texttt{dense\_\allowbreak{}down\_\allowbreak{}proj} & 65536 & 3661.14 &
2250.00 & 1631.4 & 72.5\% & 5.6\% \\
\texttt{dense\_\allowbreak{}gate\_\allowbreak{}up\_\allowbreak{}proj} &
1 & 1.00 & 8.00 & 6.1 & 0.3\% & 76.1\% \\
\texttt{dense\_\allowbreak{}gate\_\allowbreak{}up\_\allowbreak{}proj} &
8 & 7.99 & 63.90 & 50.8 & 2.3\% & 79.5\% \\
\texttt{dense\_\allowbreak{}gate\_\allowbreak{}up\_\allowbreak{}proj} &
16 & 15.95 & 127.59 & 99.4 & 4.4\% & 77.9\% \\
\texttt{dense\_\allowbreak{}gate\_\allowbreak{}up\_\allowbreak{}proj} &
64 & 63.20 & 505.57 & 384.4 & 17.1\% & 76.0\% \\
\texttt{dense\_\allowbreak{}gate\_\allowbreak{}up\_\allowbreak{}proj} &
256 & 243.61 & 1948.87 & 1219.8 & 54.2\% & 62.6\% \\
\texttt{dense\_\allowbreak{}gate\_\allowbreak{}up\_\allowbreak{}proj} &
1024 & 850.88 & 2250.00 & 1781.8 & 79.2\% & 26.2\% \\
\texttt{dense\_\allowbreak{}gate\_\allowbreak{}up\_\allowbreak{}proj} &
4096 & 2258.18 & 2250.00 & 1725.5 & 76.7\% & 9.6\% \\
\texttt{dense\_\allowbreak{}gate\_\allowbreak{}up\_\allowbreak{}proj} &
16384 & 3850.16 & 2250.00 & 1656.0 & 73.6\% & 5.4\% \\
\texttt{dense\_\allowbreak{}gate\_\allowbreak{}up\_\allowbreak{}proj} &
65536 & 4673.92 & 2250.00 & 1679.8 & 74.7\% & 4.5\% \\
\end{longtable}

\begin{answer}

\begin{enumerate}
\def\labelenumi{(\alph{enumi})}
\item
  \(M\le16\) 时约 0.1--99.4 TFLOPS；M=64/256 进入过渡。除 K=128 的
  \texttt{f\_\allowbreak{}b\_\allowbreak{}proj} 外，多数投影在
  \(M\approx1024\)--\(4096\) 接近平台，\(M\ge4096\) 的 compute 达成率约
  72.5\%--92.9\%，具体数值随形状和库 dispatch 波动。
\item
  \texttt{dense\_\allowbreak{}gate\_\allowbreak{}up\_\allowbreak{}proj}
  在 M=1/8/16 的 memory roof 达成率为 76.1\%/79.5\%/77.9\%，compute 为
  0.3\%/2.3\%/4.4\%；\texttt{in\_\allowbreak{}proj\_\allowbreak{}qkvgfab}
  的 memory roof 达成率约为
  67\%--72\%，属于带宽敏感区。这里的有效带宽按题设输入/输出字节数除以执行时间计算，缓存复用会影响该数值。
\item
  \texttt{f\_\allowbreak{}b\_\allowbreak{}proj} 的 K=128，\(M\to\infty\)
  时 AI 趋近 \(NK/(N+K)=118.15\)，低于机器平衡点 281.25。小 M 时约 4
  \(\mu\)s 的启动/调度成本占主导，MMA 初始化与少量 K
  迭代的开销占比较高，因此两个 roof 达成率都低。M=65536 时 memory
  达成率升至 73.6\%，表现为带宽敏感。
\item
  小 M 的输出工作量较少，Tensor Core 路径的
  staging、初始化和同步开销占比较高。skinny kernel 直接用 CUDA Core FMA
  读取权重，可压缩 setup 开销，这解释了题设在 \(M\le16\) 时采用该
  dispatch 的动机。
\end{enumerate}

\end{answer}

\begin{lookback}

4.1--4.3 从同一个 GEMM 出发，依次加入 tiling、TMA 和多级 pipeline。
回顾梯子表时，重点不是最终的 cuBLAS 达成率，而是能够说明每一步优化
减少了什么开销，以及新的瓶颈出现在哪里。

assignment01 Bonus 中 naive matmul 与 cuBLAS 之间的部分差距，现在已经
可以从 Tensor Core、异步数据搬运和软件流水三个方面解释。进一步的 warp
specialization、更大的 tile、epilogue 融合和 persistent kernel
等优化不在本作业要求范围内。

4.5 则说明 Tensor Core 的收益同样依赖矩阵形状。当 M 很小时，即使 kernel
本身使用了前面的优化方法，也不一定能够有效利用 Tensor Core。

\end{lookback}

\section{低精度与 block
scaling}\label{ux4f4eux7cbeux5ea6ux4e0e-block-scaling}

降低数值精度可以换取更高的 Tensor Core
吞吐，但代价是可表示的动态范围变小。本模块先从 per-tensor scale 遇到
outlier 时的问题入手，再分析 block scaling 的代数约束，最后完成一条
NVFP4 量化通路，并通过融合实验观察低精度 kernel 的实际性能。

\begin{reading}

课件 S094--S111（fp8 格式：S096；per-tensor 与 outlier：S097--S099；
block scaling 约束：S100--S101；fp4 与 NVFP4：S106--S108；SF 布局：
S108）；\texttt{cuda\_\allowbreak{}fp4.\allowbreak{}h} 中的
\texttt{\_\allowbreak{}\_\allowbreak{}nv\_\allowbreak{}fp4x2\_\allowbreak{}e2m1}；cuBLASLt
的 block-scaled FP4 matmul。格式约定见题面材料
\texttt{cuda/\allowbreak{}m5\_\allowbreak{}lowprec/\allowbreak{}nvfp4\_\allowbreak{}common.\allowbreak{}h}。

\end{reading}

\prob{5.1}{EXPERIMENT}[kernels/quant\_outlier.py]

观察 outlier 对 per-tensor scale 的影响。输入包含一万个 \([-1,1]\)
均匀分布的元素，并额外加入一个值为 3000 的 outlier。补全脚本中的两个
TODO，完成 per-tensor E4M3 量化与反量化：

\begin{verbatim}
cd assignment02
uv run python kernels/quant_outlier.py
\end{verbatim}

\begin{longtable}[]{@{}llllll@{}}
\toprule\noalign{}
采样点 \(x\approx\) & 0.5 & 0.1 & 0.01 & 0.005 & 3000 \\
\midrule\noalign{}
\endhead
\bottomrule\noalign{}
\endlastfoot
相对误差 & & & & & \\
\end{longtable}

根据实验结果回答：

\begin{enumerate}
\def\labelenumi{(\alph{enumi})}
\item
  去掉 outlier 后重新量化，\(x\approx0.5\) 处的误差变化多少倍？
\item
  找出输入被量化为 0 的阈值，并写出该阈值与 scale 的关系式。
\item
  改用 1×128 的 per-block scale 后，包含 outlier 的 block 与不包含
  outlier 的 block 分别有什么变化？
\end{enumerate}

\begin{answer}

硬件：本地 AMD Ryzen AI 9 HX 370（CPU 实验，compute capability
不适用）；nvcc 12.9.86。

本地 CPU PyTorch 2.14.0+cpu，命令
\texttt{/\allowbreak{}tmp/\allowbreak{}assignment02-\allowbreak{}venv/\allowbreak{}bin/\allowbreak{}python\ kernels/\allowbreak{}quant\_\allowbreak{}outlier.\allowbreak{}py}，原始输出
\texttt{results/\allowbreak{}local/\allowbreak{}outlier.\allowbreak{}log}。

\begin{longtable}[]{@{}lrrrrr@{}}
\toprule\noalign{}
\(x\approx\) & 0.5 & 0.1 & 0.01 & 0.005 & 3000 \\
\midrule\noalign{}
\endhead
\bottomrule\noalign{}
\endlastfoot
相对误差 & 0.04611 & 0.04634 & 0.3085 & 1.000 & 0 \\
\end{longtable}

\begin{enumerate}
\def\labelenumi{(\alph{enumi})}
\item
  去掉 outlier 后，最接近 0.5 的采样点相对误差降为
  \texttt{0.\allowbreak{}000308632181}，对应改善
  \texttt{149.\allowbreak{}40993} 倍。
\item
  E4M3 最小正 subnormal 为 \(2^{-9}\)，RN-even 舍入到零的边界为
  \(|x|\le scale\,2^{-10}\)。本次 scale=3000/448=6.69642857，理论边界
  0.00653948103；实际采样中被归零的最大绝对值为 0.00648248196。
\item
  128 元素分块后，普通 block 的平均相对误差从 0.0379851 降到
  0.0220541，归零数从 79 降到 0；outlier 所在最后一个 block
  的普通元素平均误差为 0.0257169。该 block 保留较大的 scale，其他 block
  使用各自的局部 scale；代码按实际元素数处理尾块。
\end{enumerate}

\end{answer}

\prob{5.2}{DERIVE}[kernels/block\_scale\_sim.py]

block scaling 的代数关键是沿着哪个方向分段 scale 乘积在 K
归约中才能保持常数。补全两个 fp64 模拟函数:

\begin{itemize}
\tightlist
\item
  \texttt{gemm\_\allowbreak{}scale\_\allowbreak{}per\_\allowbreak{}row\_\allowbreak{}col}:A
  每行一个 scale，B 每个输出列一个 scale。scale
  乘积在整个点积中不变，完整归约后只用乘回一次。
\item
  \texttt{gemm\_\allowbreak{}scale\_\allowbreak{}along\_\allowbreak{}k}:A
  和 B 的 scale 每 128 个 K 元素改变 一次。每个 K block 的 partial sum
  要分别乘回该段的 scale 乘积， 然后进行累加。
\end{itemize}

文件还提供了错误范例
\texttt{gemm\_\allowbreak{}scale\_\allowbreak{}along\_\allowbreak{}k\_\allowbreak{}one\_\allowbreak{}restore}：它先把
所有归一化 partial sum 相加，最后只乘回第一段的 scale。

判测:

\begin{verbatim}
cd assignment02
uv run pytest tests/test_block_scale.py
\end{verbatim}

两个正确函数只要在 fp64 容差内与直接 GEMM 等价，不要要求
bit-exact（因为分段会改变浮点加法的分组）。然后回答:

\begin{enumerate}
\def\labelenumi{(\alph{enumi})}
\item
  用两行代数式分别写出 row/column scale 为什么可以在整个 点积外乘回，而
  K-block scale 为什么必须逐段乘回。
\item
  \href{https://docs.nvidia.com/cutlass/latest/media/docs/cpp/blackwell_functionality.html}{NVIDIA
  CUTLASS 的 Blackwell SM100 GEMM 说明} 把 A 的 scale 布局写成
  \(M \times \lceil K/SV \rceil\)，B 写成
  \(N \times \lceil K/SV \rceil\)，每个 scale 负责连续的 16 或 32 个 K
  元素。结合 GEMM 内循环、连续供数和 Tensor Core 指令语义，说明硬件
  为什么把 block scale 与 K 归约对齐。
\item
  DeepSeek-V3 的 weight 128×128 表示 scale 还在输出通道方向上
  共享，但每个组仍覆盖一段 K；NVFP4 则每 16 个 K 元素一组。 结合 5.1(c)
  的误差，说明粒度 16 相对粒度 128 有什么优势，又增加了 多少 scale
  metadata 与供数复杂度。
\end{enumerate}

\begin{answer}

硬件：本地 AMD Ryzen AI 9 HX 370（CPU 实验，compute capability
不适用）；nvcc 12.9.86。

本地运行
\texttt{/\allowbreak{}tmp/\allowbreak{}assignment02-\allowbreak{}venv/\allowbreak{}bin/\allowbreak{}python\ -\allowbreak{}m\ pytest\ tests/\allowbreak{}test\_\allowbreak{}block\_\allowbreak{}scale.\allowbreak{}py\ -\allowbreak{}q}：\texttt{3\ passed}，含``错误的一次乘回''反例测试；见
\texttt{results/\allowbreak{}local/\allowbreak{}pytest.\allowbreak{}log}。CPU
PyTorch 2.14.0+cpu，运算使用 fp64。

\begin{enumerate}
\def\labelenumi{(\alph{enumi})}
\tightlist
\item
  A、B 分别按行存放，B 的行是输出列：
\end{enumerate}

\[C_{mn}=s^A_m s^B_n\sum_k(A_{mk}/s^A_m)(B_{nk}/s^B_n),\]
\[C_{mn}=\sum_j s^A_{mj}s^B_{nj}\sum_{k\in j}(A_{mk}/s^A_{mj})(B_{nk}/s^B_{nj}).\]

第一式的 scale 乘积在整个 k 归约中保持常数，可在归约后恢复；第二式的
scale 随 K block
改变，需要逐段恢复后相加。分段会改变浮点加法结合顺序，判测采用 fp64
容差比较结果。

\begin{enumerate}
\def\labelenumi{(\alph{enumi})}
\setcounter{enumi}{1}
\item
  K 是 MMA 内积归约维，一条指令连续消费一段 K；scale
  与该段绑定，使硬件可以在正确的 partial sum 上应用乘积。A 的元数据为
  \(M\times\lceil K/SV\rceil\)，B 为 \(N\times\lceil K/SV\rceil\)，既与
  tile 迭代对齐，也允许连续加载 data 和相应 scale。
\item
  16 元素组将 outlier 的影响限制在更小范围。每 scale 为 1 byte
  时，元数据数量是 128 元素组的 8 倍，即 1/16 对 1/128 B/elem。NVFP4 的
  1-byte SF 相对 0.5 B/elem 数据占 12.5\%，还需 swizzled SF 布局与指令的
  data/scale 同步供数。DeepSeek-V3 的 128×128 weight block
  同时沿输出通道方向共享 scale，其总 metadata
  比例需要结合两个方向的粒度计算。
\end{enumerate}

出处：\href{https://docs.nvidia.com/cutlass/latest/media/docs/cpp/blackwell_functionality.html}{CUTLASS
Blackwell GEMM}。

\end{answer}

\prob{5.3}{FROM-SCRATCH}[cuda/m5\_lowprec/]

完成一条 NVFP4 量化通路。开始前先阅读
\texttt{nvfp4\_\allowbreak{}common.\allowbreak{}h} 中给出的格式约定：

\begin{itemize}
\tightlist
\item
  每个量化组包含 K 方向连续的 16 个元素；
\item
  SF 使用 \texttt{amax\ /\allowbreak{}\ 6}，并以 e4m3 保存；
\item
  SF 按 Tensor Core 消费要求的 swizzled 布局存放，字节偏移公式已经给出。
\end{itemize}

\paragraph{(a) E2M1 编码}\label{a-e2m1-ux7f16ux7801}

在 \texttt{e2m1\_\allowbreak{}encode.\allowbreak{}h} 中实现 host/device
通用的 E2M1 round-to-nearest-even 编码器。

判测会在 GPU 上使用 \texttt{cuda\_\allowbreak{}fp4.\allowbreak{}h} 中的
\texttt{\_\allowbreak{}\_\allowbreak{}nv\_\allowbreak{}fp4x2\_\allowbreak{}e2m1}
转换同一批候选值，包括所有中点、边界以及采样值，并与自己的编码结果
逐位比较：

\begin{verbatim}
cd assignment02/cuda
make run/m5_lowprec/03a_encode_check
\end{verbatim}

\paragraph{(b) NVFP4 quant kernel}\label{b-nvfp4-quant-kernel}

在
\texttt{nvfp4\_\allowbreak{}quant\_\allowbreak{}kernel.\allowbreak{}h}
中实现 quant kernel，完成

\texttt{bf16\ →\ e2m1\ 打包\ +\ e4m3\ SF\ +\ swizzled\ SF\ 布局}

设备侧的 E2M1 转换直接使用
\texttt{\_\allowbreak{}\_\allowbreak{}nv\_\allowbreak{}fp4x2\_\allowbreak{}e2m1}。

第一层判测要求输出逐 byte 与 host 参考结果相等；host 参考使用你在 5.3(a)
中实现的 E2M1 编码器生成。随后将量化结果原样交给 cuBLASLt 的 FP4
matmul，检查生成的数据和 SF 布局能否被实际消费：

\begin{verbatim}
make run/m5_lowprec/03b_nvfp4_quant
make run/m5_lowprec/test_fp4_gemm
\end{verbatim}

正确实现下，FP4 matmul 的 \texttt{maxrel} 约为
\texttt{4e-\allowbreak{}3}，主要来自 bf16 输出的舍入误差。如果 SF
布局错位或 scale 对应到了错误的量化组，
通常会出现成块的明显数值错误，而不是仅有小幅舍入误差。

\paragraph{(c) Ceiling probe}\label{c-ceiling-probe}

在
\texttt{03c\_\allowbreak{}ceiling\_\allowbreak{}probe.\allowbreak{}cu}
中实现一个 ceiling probe。它与 quant kernel
使用相同的访存模式，但不执行量化计算，只读取数据并通过简单 xor 后写回。

报告：

\begin{itemize}
\tightlist
\item
  ceiling probe 的 GB/s；
\item
  quant kernel 的 GB/s；
\item
  两者的比值。
\end{itemize}

结合 Nsight Compute 判断 quant kernel 距离自己的访存上限还有多远，
以及剩余差距主要来自访存还是计算。

\paragraph{(d) Optional}\label{d-optional}

使用 tcgen05 \texttt{kind:\allowbreak{}:\allowbreak{}mxf4nvf4}
直接消费自己生成的 NVFP4 数据， SF 通过 TMEM 提供，并与 cuBLASLt
的结果对拍。

\begin{answer}

硬件：NVIDIA B300 SXM6 AC，compute capability 10.3，148 SM；CUDA 13.0 /
nvcc 13.0.88，\texttt{ARCH=\allowbreak{}100f}。

\begin{enumerate}
\def\labelenumi{(\alph{enumi})}
\item
  \texttt{e2m1\_\allowbreak{}encode}
  保留符号（含负零），在幅值中点处选择编码低位为偶数的一侧，并将大值饱和到
  6：0.25→0、0.75→1、1.25→1、1.75→2、2.5→2、3.5→4、5→4。\texttt{make\ -\allowbreak{}B\ run/\allowbreak{}m5\_\allowbreak{}lowprec/\allowbreak{}03a\_\allowbreak{}encode\_\allowbreak{}check}：\texttt{PASS:\allowbreak{}\ 202864\ values\ match\ hardware}。
\item
  一个线程负责连续 16 个 bf16，用两次 uint4 读取，计算 amax→E4M3
  SF→反量化 SF 的倒数，然后硬件
  \texttt{\_\allowbreak{}\_\allowbreak{}nv\_\allowbreak{}fp4x2\_\allowbreak{}e2m1}
  编码，八个 byte 合并为 uint64 写回。SF 使用题面 swizzled
  offset，padding 由调用者清零。
\end{enumerate}

\begin{Shaded}
\begin{Highlighting}[]
\FunctionTok{make} \AttributeTok{{-}B}\NormalTok{ run/m5\_lowprec/03b\_nvfp4\_quant}
\FunctionTok{make} \AttributeTok{{-}B}\NormalTok{ run/m5\_lowprec/test\_fp4\_gemm}
\end{Highlighting}
\end{Shaded}

quant 的 128×1024、200×4096、4096×7168 都
\texttt{PASS(bad=\allowbreak{}0)}，同时检查 data/SF bytes。cuBLASLt
三个形状均 PASS，maxrel 分别为 0.003880、0.003891、0.003880，说明打包与
scale 布局可直接消费。

\begin{enumerate}
\def\labelenumi{(\alph{enumi})}
\setcounter{enumi}{2}
\tightlist
\item
  probe 与 quant 使用相同 grid、block、两次 uint4 输入和 8B data+1B SF
  输出。探针对所有输入 word 做 XOR
  后写回，使读入数据参与输出计算。有效字节为 \(2+0.5+1/16=2.5625\)
  B/elem。
\end{enumerate}

题面原始主程序在 4096×7168 上测得 probe=7301 GB/s、quant=5425 GB/s，比值
0.743；同一进程追加对照（\texttt{experiments/\allowbreak{}bandwidth.\allowbreak{}cu}）为
7323.6/5590.5 GB/s，比值 0.7634。以上性能数据采用独立计时，profiler
数据用于归因。

大形状 NCU 的 quant 为 SM throughput 51.56\%、DRAM
28.99\%，有效带宽达到同形探针约 76\%。结合这组指标，剩余开销包括
amax、转换、倒数、packing 与指令发射。其它形状见
\texttt{bandwidth.\allowbreak{}log}；profiler 记录见
\texttt{ncu-\allowbreak{}quant-\allowbreak{}large.\allowbreak{}csv}。

\begin{enumerate}
\def\labelenumi{(\alph{enumi})}
\setcounter{enumi}{3}
\tightlist
\item
  未选做手写
  \texttt{tcgen05.\allowbreak{}kind:\allowbreak{}:\allowbreak{}mxf4nvf4}，已完成要求的
  cuBLASLt 消费验证。
\end{enumerate}

\end{answer}

\begin{lookback}

5.3(c) 用来区分量化 kernel 中的访存成本和数值转换成本。作为参考，
软件实现 E2M1 RN-even 时，每个 kernel 约执行 18129 条 FSETP；
使用硬件转换后约为 1272 条 \texttt{F2FP.\allowbreak{}E2M1}，同一 kernel
的时间从 71.7 µs 降到 32.8 µs。Nsight Compute 中的瓶颈也从 SM 利用率
84.7\% 的 compute-bound 转为 memory-bound。

这组结果说明，当低精度转换由专用硬件完成后，量化 kernel 可以更接近
纯数据搬运的性能特征；5.3(c) 的 ceiling probe 就用于衡量这一差距。

Hopper 上的 fp8 累加需要软件 promotion（S102--S105，DeepGEMM 使用
双层累加循环），而 Blackwell 进一步在硬件中支持 block scaling
（S105）。本作业不单独设置对应题目，但 5.2 中的 scale 分组约束是
理解两者的共同基础。

\end{lookback}

\begin{capstone}{prob 5.4(FROM-SCRATCH):融合 rms\_norm + NVFP4}

实现融合的 \texttt{rms\_\allowbreak{}norm\ +\ NVFP4} kernel：

\[
y = \mathrm{rms\_norm}(x)\cdot w
\]

要求将结果直接量化为 NVFP4，不写回 bf16 中间结果。

本题以 vLLM fused kernel 清单中的 issue \#25179，以及两个相关实现 PR
\#36413、\#32957 为背景。相关实现的正确性没有问题，也确实将 kernel
数量从 2 个减少到了 1 个，但当时观察到的端到端收益仍处于 ±1\% 左右的
测量噪声内。本题要通过逐形状实验解释这部分收益去了哪里。

按理论字节量计算，两步实现约为 6.56 B/elem，融合后约为 2.56
B/elem，因此单纯从访存量估计可以得到约 2.56× 的理想加速。
这个数字只是上限，实际结果需要通过实验验证。

程序已经提供两步基线、正确性检查和逐形状计时框架。测试覆盖：

\begin{itemize}
\tightlist
\item
  \(M\) 从 1 到 16384；
\item
  \(K \in \{4096, 7168, 8192\}\)；
\item
  共十个测试形状。
\end{itemize}

需要完成三部分：

\begin{enumerate}
\def\labelenumi{\arabic{enumi}.}
\tightlist
\item
  实现融合 kernel，具体线程组织和 tile 结构不限；
\item
  分别调优融合实现和两步基线，使两边都使用各自合理的配置后再比较；
\item
  完成逐形状性能分析。
\end{enumerate}

第二点必须认真处理。让两步基线处于明显不利的配置下得到的加速比没有
比较意义，这也是相关上游 PR 的实验过程中暴露出的一个问题。

运行：

\begin{verbatim}
cd assignment02/cuda
make run/m5_lowprec/04_fused_rms_nvfp4
\end{verbatim}

报告逐形状结果，包括：

\begin{itemize}
\tightlist
\item
  实测加速比；
\item
  融合 kernel 相对 5.3(c) ceiling probe 的性能比例。
\end{itemize}

重点解释实测结果与 2.56× 理论上限之间的差距：不同 M 范围内主要受到
什么因素限制，并给出 Nsight Compute 数据或计算结果作为依据。

本题不设置性能门槛。实现首先需要保证正确，评分重点放在实验设计和性能归因。

Optional：根据分析得到的主要瓶颈进行一次针对性优化，重新测试并更新表格。

\end{capstone}

\begin{answer}

硬件：NVIDIA B300 SXM6 AC，compute capability 10.3，148 SM；CUDA 13.0 /
nvcc 13.0.88，\texttt{ARCH=\allowbreak{}100f}。

实现 \texttt{rms\_\allowbreak{}quant\textless{}B\textgreater{}}：CTA
内归约 sumsq，输入以 uint4 向量读入并保留在寄存器中，计算 rnorm
后直接将归一化结果量化成 FP4 并写回。对 K=4096/7168/8192
按实际组长度处理尾部。

\begin{Shaded}
\begin{Highlighting}[]
\FunctionTok{make} \AttributeTok{{-}B}\NormalTok{ run/m5\_lowprec/04\_fused\_rms\_nvfp4}
\end{Highlighting}
\end{Shaded}

每个形状分别扫描融合
\texttt{B=\allowbreak{}128/\allowbreak{}256/\allowbreak{}512}，基线 RMS
\texttt{B=\allowbreak{}128/\allowbreak{}256/\allowbreak{}512} 与
\texttt{grid=\allowbreak{}M/\allowbreak{}2SM/\allowbreak{}4SM/\allowbreak{}8SM}
的组合，选定配置后独立计时。quant 沿用已验证的 uint4/128-thread
grid-stride 实现。完整候选时间与所选配置保存在
\texttt{results/\allowbreak{}b300/\allowbreak{}fused-\allowbreak{}tuned.\allowbreak{}log}。

\begin{longtable}[]{@{}
  >{\raggedright\arraybackslash}p{(\linewidth - 10\tabcolsep) * \real{0.1304}}
  >{\raggedleft\arraybackslash}p{(\linewidth - 10\tabcolsep) * \real{0.1739}}
  >{\raggedleft\arraybackslash}p{(\linewidth - 10\tabcolsep) * \real{0.1739}}
  >{\raggedleft\arraybackslash}p{(\linewidth - 10\tabcolsep) * \real{0.1739}}
  >{\raggedleft\arraybackslash}p{(\linewidth - 10\tabcolsep) * \real{0.1739}}
  >{\raggedleft\arraybackslash}p{(\linewidth - 10\tabcolsep) * \real{0.1739}}@{}}
\toprule\noalign{}
\begin{minipage}[b]{\linewidth}\raggedright
M×K
\end{minipage} & \begin{minipage}[b]{\linewidth}\raggedleft
两步 \(\mu\)s
\end{minipage} & \begin{minipage}[b]{\linewidth}\raggedleft
融合 \(\mu\)s
\end{minipage} & \begin{minipage}[b]{\linewidth}\raggedleft
加速比
\end{minipage} & \begin{minipage}[b]{\linewidth}\raggedleft
融合/同形 ceiling 带宽
\end{minipage} & \begin{minipage}[b]{\linewidth}\raggedleft
data/SF 不同 byte 数
\end{minipage} \\
\midrule\noalign{}
\endhead
\bottomrule\noalign{}
\endlastfoot
1×4096 & 7.18 & 4.11 & 1.75× & 76.0\% & 0/0 \\
16×4096 & 8.20 & 4.11 & 2.00× & 74.8\% & 0/0 \\
256×4096 & 8.21 & 4.11 & 2.00× & 99.9\% & 0/0 \\
1024×4096 & 10.25 & 6.16 & 1.67× & 67.0\% & 0/0 \\
4096×4096 & 20.53 & 14.42 & 1.42× & 56.9\% & 0/0 \\
16384×4096 & 79.44 & 59.47 & 1.34× & 52.8\% & 6/1 \\
4096×7168 & 38.91 & 20.59 & 1.89× & 49.9\% & 3/0 \\
16384×7168 & 134.25 & 88.83 & 1.51× & 58.3\% & 5/1 \\
4096×8192 & 41.83 & 22.58 & 1.85× & 54.6\% & 6/2 \\
16384×8192 & 149.17 & 92.29 & 1.62× & 63.6\% & 2/1 \\
\end{longtable}

十个形状均 PASS。检查覆盖 data 和 SF（含 padding），均按题面的 1e-4
比例容差处理。少量 byte 差异来自 sumsq
加法分组使中点附近的值翻转。两步基线在中间写 bf16
时额外舍入，融合路径直接量化归一化结果；计时比较的是这两条精度路径。

比例用 \texttt{bandwidth.\allowbreak{}log} 的同形 probe
时间除以本表融合时间，两者统一按 2.5625 B/elem
计算有效带宽。该比例衡量融合实现相对 quant
访存探针的速度；融合还包含权重读取与归约。

小 M 时 kernel 启动成本主导，M=1 的 CTA 数量远少于 148
SM，融合通过减少一次启动获得主要收益。中等 M 的 CTA
数量增加，归约、SF/FP4 转换的发射与同步成本开始显现。大 M
的开销包括寄存器占用、权重读取、归约与转换；NCU 对 4096×7168
的融合实测为 52 registers/thread、active warp 43.31\%、SM throughput
52.17\%、DRAM 18.93\%。

2.56×
的理想估算按两步与融合的字节量之比计算，并假定两者以相同带宽执行。实测还包含启动、权重读取、转换和归约开销，最终加速比为
1.34--2.00×。

针对性优化：初稿逐标量读取并在归约后重读 x，大 M 加速比为
0.53--0.64×；改为向量装载、寄存器复用，并分别调优融合与基线后，得到上表结果。原始首轮见
\texttt{04\_\allowbreak{}fused\_\allowbreak{}rms\_\allowbreak{}nvfp4.\allowbreak{}log}，完整候选配置与最终计时见
\texttt{fused-\allowbreak{}tuned.\allowbreak{}log}。

\end{answer}

\prob{5.5}{CONCEPT}

比较 W4A16 + Marlin kernel（权重 int4、计算 fp16）与 5.3 中的 NVFP4
GEMM。

根据课件 S109 的分类回答：

\begin{enumerate}
\def\labelenumi{(\alph{enumi})}
\item
  两者分别属于存储量化还是计算量化？
\item
  两种方法分别节省哪些资源：显存容量、显存带宽还是计算吞吐？
\item
  在 4.5 中的小 batch decode 场景下，哪一类量化的收益更加直接？
\end{enumerate}

每问用两到三句话回答。

\begin{answer}

硬件：NVIDIA B300 SXM6 AC，compute capability 10.3，148 SM；CUDA 13.0 /
nvcc 13.0.88，\texttt{ARCH=\allowbreak{}100f}。

\begin{enumerate}
\def\labelenumi{(\alph{enumi})}
\item
  W4A16/Marlin 是权重存储量化，int4 权重解码后仍执行 fp16 计算；NVFP4
  是计算量化，原生 FP4 Tensor Core 直接消费打包 FP4 数据与 block
  scales。
\item
  两者都减少权重容量和读取流量；Marlin 的收益需扣除解码成本，Tensor Core
  精度仍为 fp16。NVFP4 还压缩 activation，并可使用更高的 FP4
  计算峰值，但需要量化、scale 元数据及供数布局。
\item
  小 batch decode 的权重搬运占比高，W4A16
  通过减少权重字节数获得直接收益，同时保持 activation 为 fp16。NVFP4
  同样节省带宽；随着计算规模增大，量化与启动开销得到摊薄，更高的 FP4
  计算吞吐逐步发挥作用。
\end{enumerate}

\end{answer}

\section{TileLang 对照}\label{tilelang-ux5bf9ux7167}

前面的模块已经手动处理过 Tensor Core 指令、数据布局和数据搬运。
本模块通过 TileLang 的 lowering 结果观察其中哪些工作可以交给 DSL 完成。

\begin{reading}

课件 S112；assignment01 Module 7（7.5 的表、7.6 的
\texttt{kernels/\allowbreak{}tilelang\_\allowbreak{}matmul.\allowbreak{}py}）；TileLang
文档中的 lowering/debug 部分。

\end{reading}

\prob{6.1}{EXPERIMENT}

使用 assignment01 的
\texttt{kernels/\allowbreak{}tilelang\_\allowbreak{}matmul.\allowbreak{}py}，或任意使用
\texttt{T.\allowbreak{}gemm} 的 kernel，分别以
\texttt{sm\_\allowbreak{}90a} 和 \texttt{sm\_\allowbreak{}100a} 为
target 编译。 本题只要求编译，不需要实际使用 H 卡。

保存生成的 CUDA 源码与 lowering 输出，并填写：

\begin{longtable}[]{@{}lll@{}}
\toprule\noalign{}
& \texttt{sm\_\allowbreak{}90a} & \texttt{sm\_\allowbreak{}100a} \\
\midrule\noalign{}
\endhead
\bottomrule\noalign{}
\endlastfoot
选中的 Tensor Core 指令 & & \\
descriptor 在哪里、由谁生成 & & \\
smem swizzle 布局在哪一步确定 & & \\
数据由谁搬入 smem & & \\
\end{longtable}

与 M2--M4 中的手写实现进行对照，并回答：

\begin{enumerate}
\def\labelenumi{(\alph{enumi})}
\item
  哪些硬件相关的决策已经由 DSL 自动完成？
\item
  哪些参数仍然需要程序员决定，例如 tile 尺寸和 stages？
\end{enumerate}

最后在 assignment01 7.5 的``谁负责''表中补充一行：

\texttt{Tensor\ Core\ 指令选择与供数布局}

\begin{answer}

硬件：NVIDIA B300 SXM6 AC，compute capability 10.3，148 SM；CUDA 13.0 /
nvcc 13.0.88，\texttt{ARCH=\allowbreak{}100f}。

固定 TileLang 0.1.13；使用
\texttt{experiments/\allowbreak{}tilelang\_\allowbreak{}lower.\allowbreak{}py}
的 T.gemm，tile=128×128×64、threads=128、stages=3，分别 lower 到
\texttt{sm\_\allowbreak{}90a/\allowbreak{}sm\_\allowbreak{}100a，再显式调用}
nvcc 生成对应 cubin。两种架构均完成编译验证。

\begin{Shaded}
\begin{Highlighting}[]
\ExtensionTok{.venv/bin/python}\NormalTok{ experiments/tilelang\_lower.py}
\end{Highlighting}
\end{Shaded}

输入 IR、device/host lowering、CUDA 源码和编译命令/返回值在
\texttt{results/\allowbreak{}tilelang/\allowbreak{}}，运行记录在
\texttt{results/\allowbreak{}b300/\allowbreak{}tilelang.\allowbreak{}log}。

\begin{longtable}[]{@{}
  >{\raggedright\arraybackslash}p{(\linewidth - 4\tabcolsep) * \real{0.3333}}
  >{\raggedright\arraybackslash}p{(\linewidth - 4\tabcolsep) * \real{0.3333}}
  >{\raggedright\arraybackslash}p{(\linewidth - 4\tabcolsep) * \real{0.3333}}@{}}
\toprule\noalign{}
\begin{minipage}[b]{\linewidth}\raggedright
项目
\end{minipage} & \begin{minipage}[b]{\linewidth}\raggedright
\texttt{sm\_\allowbreak{}90a}
\end{minipage} & \begin{minipage}[b]{\linewidth}\raggedright
\texttt{sm\_\allowbreak{}100a}
\end{minipage} \\
\midrule\noalign{}
\endhead
\bottomrule\noalign{}
\endlastfoot
此配置实际选择的指令 &
\texttt{tl:\allowbreak{}:\allowbreak{}wgmma\_\allowbreak{}ss\textless{}.\allowbreak{}.\allowbreak{}.\allowbreak{},\allowbreak{}64,\allowbreak{}128,\allowbreak{}16,\allowbreak{}.\allowbreak{}.\allowbreak{}.\allowbreak{}\textgreater{}}
→ WGMMA &
\texttt{tl:\allowbreak{}:\allowbreak{}mma\_\allowbreak{}sync\textless{}.\allowbreak{}.\allowbreak{}.\allowbreak{},\allowbreak{}16,\allowbreak{}8,\allowbreak{}16,\allowbreak{}.\allowbreak{}.\allowbreak{}.\allowbreak{}\textgreater{}}
→ MMA.SYNC \\
MMA descriptor & device 端
\texttt{initialize\_\allowbreak{}wgmma\_\allowbreak{}descriptor} 构造
A/B 的 64-bit descriptor & warp fragment 路径无 WGMMA/tcgen05 MMA
descriptor \\
TMA tensor map & host lowering 写入 encode 参数，launch
前运行库编码，kernel 参数为
\texttt{a\_\allowbreak{}desc/\allowbreak{}b\_\allowbreak{}desc} &
同左 \\
shared swizzle & 编译期 LayoutInference / LowerTileOp
确定，落实为描述符参数、shared 索引和 TMA 布局 &
同样由布局推断确定，落实为 TMA shared 布局及 fragment load 地址 \\
谁搬到 shared &
\texttt{tl:\allowbreak{}:\allowbreak{}tma\_\allowbreak{}load}
和自动生成的 mbarrier/pipeline & 同为
\texttt{tl:\allowbreak{}:\allowbreak{}tma\_\allowbreak{}load}，随后 warp
级 fragment 装载 \\
\end{longtable}

\begin{enumerate}
\def\labelenumi{(\alph{enumi})}
\item
  编译器完成指令选择、fragment 分配、shared 布局、TMA map
  参数与同步插入。在 TileLang 0.1.13
  \texttt{的本次配置中，sm\_\allowbreak{}90a} 选择
  \texttt{WGMMA，sm\_\allowbreak{}100a} 选择 MMA.SYNC，两者均由 TMA
  搬运数据。
\item
  程序员选择 tile
  \texttt{尺寸、threads、num\_\allowbreak{}stages、输入和累加}
  dtype，并通过实验验证正确性、调整性能。与手写 M2--M4 相比，DSL
  承担布局推断、指令生成和同步插入，程序员侧重计算划分与参数调优。
\end{enumerate}

已在 assignment01 7.5 的作答表中补充``Tensor Core
指令选择与供数布局''一行。此版本 target 使用 dict，并在
\texttt{with\ Target(.\allowbreak{}.\allowbreak{}.\allowbreak{})}
上下文内调用 lower；旧 CLI 风格 target 字符串会报错。

\end{answer}

\section*{团队选做（推荐）}\label{ux56e2ux961fux9009ux505aux63a8ux8350}
\addcontentsline{toc}{section}{团队选做（推荐）}

2--4 人一组，从 \texttt{team/\allowbreak{}}
中的两个题目任选一个，也可以全部完成：

\begin{itemize}
\tightlist
\item
  C1：FlashKDA 官方 kernel 当前使用 SM80 MMA，分析迁移到 SM100
  是否值得；
\item
  C2：MiniMax M3 MSA 的小 batch decode 当前使用
  Triton，分析是否值得实现专用 kernel。
\end{itemize}

每道题分为三个阶段：

\texttt{复现与测量\ →\ 分析（结论\ +\ 证据）→\ 挑战}

挑战阶段不要求一定得到正向加速。如果实验结果表明继续优化收益有限，
只要能够用数据说明为什么不值得继续做，同样视为有效结论。

答辩时间为 10 分钟，另有 5 分钟提问；提问优先由选择另一道团队题的
小组提出。具体要求见 \texttt{team/\allowbreak{}README.\allowbreak{}md}。

团队题不限制 AI 工具的使用。

\section*{提交}\label{ux63d0ux4ea4}
\addcontentsline{toc}{section}{提交}

\begin{itemize}
\tightlist
\item
  \textbf{代码}：提交所有动手题的实现与判测输出；FROM-SCRATCH
  题同时保留判测脚本的 PASS 记录。
\item
  \textbf{报告}：包含纸面题解答、实验表格与性能归因，以及 DEBUG
  题的现象记录和修改说明。所有实验数据注明使用的 GPU。
\item
  \textbf{团队题}：提交代码、报告并完成答辩，具体要求见
  \texttt{team/\allowbreak{}README.\allowbreak{}md}。
\end{itemize}

\end{document}
